\subsection{Isotypical Components of a Module}

DEFINITION 4. --- Let A be a ring, M an A-module, and S a simple A-module. The isotypical component of type S of M, denoted by \(M_{S}\) , is the sum of the submodules of M isomorphic to S.

It is clear that \(M_{S}\) is the greatest submodule of M that is isotypical of type S. Because every submodule of \(M_{S}\) is isotypical of type S (VIII, p.~61, Proposition 2), we have \(N_{S} = M_{S} \cap N\) for every submodule N of M.

If \(S'\) is a simple A-module isomorphic to S, then we clearly have \(M_{S} = M_{S'}\) , so \(M_{S}\) depends only on the class of S (VIII, p.~51).

Let M be an A-module. There exists a greatest semisimple submodule of M, called the socle of M; it is the sum of the simple submodules of M and also the sum of the isotypical components of M. In particular, M is semisimple if and only if it is equal to its socle.

PROPOSITION 4. --- Let A be a ring. Denote the set of classes of simple A-modules by S. Let M be a semisimple A-module.

\begin{enumerate}
\def\labelenumi{\alph{enumi})}
\item
  The module M is the direct sum of the family \((\mathrm{M}_{\lambda})_{\lambda\in\mathcal{S}}\) of its isotypical components.
\item
  Suppose that M is the direct sum of a family \((\mathrm{N}_{i})_{i\in\mathrm{I}}\) of simple submodules. For every \(\lambda\in\mathscr{S}\) , let \(\mathrm{I}(\lambda)\) be the set of indices \(i\in I\) such that \(N_{i}\) is of class \(\lambda\) . We have \(M_{\lambda}=\bigoplus_{i\in\mathrm{I}(\lambda)}N_{i}\) .
\item
  If M is finitely generated, then the set of \(\lambda\in S\) such that \(M_{\lambda}\neq0\) is finite.
\item
  For every submodule \(\mathrm{N}\) of \(\mathrm{M}\) and every \(\lambda \in \mathcal{S}\) , we have \(\mathrm{N}_{\lambda} = \mathrm{N} \cap \mathrm{M}_{\lambda}\) and \((\mathrm{M} / \mathrm{N})_{\lambda} = (\mathrm{M}_{\lambda} + \mathrm{N}) / \mathrm{N}\) .
\end{enumerate}

Since M is semisimple, it is the sum of the family \((\mathrm{M}_{\lambda})_{\lambda\in\mathcal{S}}\) ; let us prove that this sum is direct. Let \(\lambda\in S\) . Denote the sum of the family \((\mathrm{M}_{\mu})_{\mu\in\mathcal{S}-\{\lambda\}}\) by \(M_{\lambda}^{\prime}\) . The module \(M_{\lambda}^{\prime}\) is the direct sum of a family of simple modules not isomorphic to \(\lambda\) (VIII, p.~56, Theorem 1). By Corollary 3 of Theorem 1 of VIII, p.~56, \(M_{\lambda}^{\prime}\) does not contain any simple submodules of class \(\lambda\) . We consequently have \(M_{\lambda}\cap M_{\lambda}^{\prime}=0\) . Assertion a) is therefore proved. By construction, we have \(M_{\lambda}\supset\bigoplus_{i\in\mathrm{I}(\lambda)}N_{i}\) , so assertion b) follows from Remark 1 of II, §1, No.~8, p.~208.

Assertion c) follows from a) and Proposition 23 of II, §1, No.~12, p.~221.

Let N be a submodule of M and \(\lambda\in S\) . The isotypical component \(N_{\lambda}\) of type \(\lambda\) of N is contained in \(M_{\lambda}\) and \(M_{\lambda}\cap N\subset N_{\lambda}\) . The intersection \(N\cap M_{\lambda}\) is therefore the isotypical component of N of type \(\lambda\) .

For every \(\lambda \in \mathcal{S}\) , the module \(\mathrm{M}_{\lambda} + \mathrm{N} / \mathrm{N}\) is isomorphic to \(\mathrm{M}_{\lambda} / (\mathrm{M}_{\lambda} \cap \mathrm{N})\) . It is therefore isotypical of type \(\lambda\) and contained in \((\mathrm{M} / \mathrm{N})_{\lambda}\) . The last assertion then follows from a) and II, §1, No.~8, p.~208, Remark 1.

COROLLARY. --- Let A be a ring and S the set of classes of simple A-modules. Let M be a semisimple A-module and N a submodule of M. Then we have \(N = \bigoplus_{\lambda \in \mathscr{S}} N \cap M_{\lambda}\) and \(M/N = \bigoplus_{\lambda \in \mathscr{S}} (M_{\lambda} + N)/N\) .

Since N and M/N are semisimple (VIII, p.~56, Corollary 3), the corollary follows from Proposition 4, d).

The support of a semisimple A-module M is the set of classes \(\lambda\) of simple A-modules such that the isotypical component of M of type \(\lambda\) is nonzero. The support of a finitely generated semisimple A-module is finite.

PROPOSITION 5. --- Let A be a ring, and let S be the set of classes of simple A-modules. Let M and N be A-modules.

\begin{enumerate}
\def\labelenumi{\alph{enumi})}
\item
  Let \(f: \mathrm{M} \to \mathrm{N}\) be a homomorphism. For every \(\lambda \in \mathcal{S}\) , \(f\) induces a homomorphism \(f_{\lambda}\) from \(\mathrm{M}_{\lambda}\) to \(\mathrm{N}_{\lambda}\) ; if \(\mathrm{M}\) is semisimple and \(f\) surjective, then each of the homomorphisms \(f_{\lambda}\) is surjective.
\item
  Suppose that M is semisimple. The mapping \(f \mapsto (f_{\lambda})_{\lambda \in \mathcal{S}}\) is a group isomorphism from \(\mathrm{Hom}_{\mathrm{A}}(\mathrm{M},\mathrm{N})\) to \(\prod_{\lambda \in \mathcal{S}}\mathrm{Hom}_{\mathrm{A}}(\mathrm{M}_{\lambda},\mathrm{N}_{\lambda})\) . When M is equal to N, the mapping is a ring isomorphism from \(\mathrm{End}_{\mathrm{A}}(\mathrm{M})\) to \(\prod_{\lambda \in \mathcal{S}}\mathrm{End}_{\mathrm{A}}(\mathrm{M}_{\lambda})\) .
\end{enumerate}

For every \(\lambda \in \mathcal{S}\) , the submodule \(f(\mathrm{M}_{\lambda})\) of \(\mathbf{N}\) is isomorphic to a quotient of an isotypical module of type \(\lambda\) ; it is therefore isotypical of type \(\lambda\) and, consequently, contained in \(\mathrm{N}_{\lambda}\) .

Suppose that \(\mathrm{M}\) is semisimple and \(f\) surjective. Then \(f\) induces an isomorphism from \(\mathrm{M} / \mathrm{Ker}(f)\) to \(\mathrm{N}\) that sends \((\mathrm{M}_{\lambda} + \mathrm{Ker}(f)) / \mathrm{Ker}(f)\) to \(f(\mathrm{M}_{\lambda})\) .

By Proposition 4 of VIII, p.~65, we have \(N_{\lambda} = f(M_{\lambda})\) , which completes the proof of a).

The mapping considered in b) is clearly a group homomorphism, and it is a ring homomorphism when M is equal to N. Let \((f_{\lambda})_{\lambda \in \mathcal{S}}\) be an element of \(\prod_{\lambda \in \mathcal{S}} \mathrm{Hom}(M_{\lambda}, N_{\lambda})\) . The unique element of its inverse image under the mapping in b) is the homomorphism \(f : M \to N\) defined by

\[
f \Bigl (\sum_ {\lambda \in \mathcal {S}} x _ {\lambda} \Bigr) = \sum_ {\lambda \in \mathcal {S}} f _ {\lambda} (x _ {\lambda})
\]

for every \((x_{\lambda})_{\lambda \in \mathcal{S}}\in \bigoplus_{\lambda}\mathrm{M}_{\lambda}\).

Remark. --- Let A and B be rings. Let M be an (A, B)-bimodule. It follows from Proposition 5 that the isotypical components of the A-module M are sub-bimodules of M. This holds, in particular, when M is an A-module and B is the opposite ring of \(\operatorname{End}_{\mathrm{A}}(\mathrm{M})\) .

Example. --- Let us consider the case when the ring A is commutative. The mapping that sends a maximal ideal m onto \(\operatorname{cl}(A/m)\) is a bijection from the set of maximal ideals of A to the set S of classes of simple A-modules (VIII, p.~51). The inverse bijection sends \(\lambda\) to its annihilator \(m_{\lambda}\) .

Let M be an A-module. For every \(\lambda\in S\) , the isotypical component \(M_{\lambda}\) of type \(\lambda\) of M consists of the elements annihilated by \(m_{\lambda}\) , and we can view \(M_{\lambda}\) as a vector space over the field \(A/m_{\lambda}\) . If M is semisimple and N is another A-module, then we can deduce a group isomorphism from \(\operatorname{Hom}_{\mathrm{A}}(\mathrm{M},\mathrm{N})\) to \(\prod_{\lambda\in\mathcal{S}}\operatorname{Hom}_{\mathrm{A}/m_{\lambda}}(\mathrm{M}_{\lambda},\mathrm{N}_{\lambda})\) from Proposition 5.

COROLLARY 1. --- Let M be a semisimple A-module and N a submodule of M. The following properties are equivalent:

\begin{enumerate}
\def\labelenumi{(\roman{enumi})}
\item
  There exists a unique submodule supplementary to N in M.
\item
  We have \(\mathrm{Hom}_{\mathrm{A}}(\mathrm{M} / \mathrm{N},\mathrm{N}) = 0\)
\item
  There exists a subset \(\Lambda\) of \(\mathcal{S}\) such that \(N = \bigoplus_{\lambda \in \Lambda} M_{\lambda}\) .
\end{enumerate}

Choose a submodule \(N^{\prime}\) supplementary to N in M (VIII, p.~56, Corollary 2). If we identify M with \(N^{\prime} \times N\) , then the submodules of M supplementary to N are the graphs of the A-linear mappings from \(N^{\prime}\) to N. Since \(N^{\prime}\) is isomorphic to M/N, we have proved the equivalence of properties (i) and (ii). By Proposition 5, b), the group \(\operatorname{Hom}_{\mathrm{A}}(\mathrm{N}^{\prime}, \mathrm{N})\) is isomorphic to the group \(\prod_{\lambda \in \mathcal{S}} \operatorname{Hom}_{\mathrm{A}}(\mathrm{N}_{\lambda}^{\prime}, \mathrm{N}_{\lambda})\) . It is zero if and only if for every \(\lambda \in S\) , we have \(N_{\lambda} = 0\) or \(N_{\lambda}^{\prime} = 0\) (VIII, p.~61, Remark), that is, \(N_{\lambda} = 0\) or \(N_{\lambda} = M_{\lambda}\) . This proves the equivalence of properties (ii) and (iii).

COROLLARY 2. --- Let M be an A-module. The following two conditions are equivalent:

\begin{enumerate}
\def\labelenumi{(\roman{enumi})}
\item
  Every submodule of M admits a unique supplementary submodule.
\item
  The module M is the direct sum of a family \((\mathrm{S}_{i})_{i\in\mathrm{I}}\) of simple, pairwise nonisomorphic modules.
\end{enumerate}

Suppose that M satisfies these conditions. Then, for every submodule N of M, there exists a unique subset J of I such that we have \(N = \bigoplus_{j \in J} S_j\) and every simple submodule of M is equal to one of the \(S_i\) .

Conditions (i) and (ii) both imply that M is semisimple.

Suppose that condition (i) is satisfied. Let \(\lambda \in \mathcal{S}\) . By the equivalence of (i) and (iii) in Corollary 1, every submodule of \(\mathrm{M}_{\lambda}\) is zero or equal to \(\mathrm{M}_{\lambda}\) . Consequently, \(\mathrm{M}_{\lambda}\) is zero or simple, and (ii) follows from the fact that \(\mathrm{M}\) is the direct sum of the family \((\mathrm{M}_{\lambda})_{\lambda \in \mathcal{S}}\) .

Conversely, if condition (ii) is satisfied, then \(M_{\lambda}\) is zero or simple for every \(\lambda \in S\) . If N is a submodule of M, then \(N \cap M_{\lambda}\) is equal to 0 or \(M_{\lambda}\) for every \(\lambda \in S\) . Since we have \(N = \bigoplus_{\lambda \in S} (N \cap M_{\lambda})\) (VIII, p.~66, Corollary), the submodule N has property (iii) of Corollary 1 and therefore admits a unique supplementary submodule in M. This proves that (ii) implies (i), as well as the last assertions of the corollary.

Let M be an A-module and S a simple A-module. Denote the opposite ring of the field \(\operatorname{End}_{\mathrm{A}}(\mathrm{S})\) by D, and view S as an (A,D)-bimodule. Then \(\operatorname{Hom}_{\mathrm{A}}(\mathrm{S},\mathrm{M})\) is a left vector space over D, and \(\operatorname{Hom}_{\mathrm{A}}(\mathrm{M},\mathrm{S})\) is a right vector space over D. The dual of the left D-vector space \(\operatorname{Hom}_{\mathrm{A}}(\mathrm{S},\mathrm{M})\) is a right vector space over D (II, §2, No.~3, p.~232, Definition 2). For every \(u \in \operatorname{Hom}_{\mathrm{A}}(\mathrm{M},\mathrm{S})\) , the mapping \(h(u): v \mapsto u \circ v\) from \(\operatorname{Hom}_{\mathrm{A}}(\mathrm{S},\mathrm{M})\) to \(\operatorname{Hom}_{\mathrm{A}}(\mathrm{S},\mathrm{S}) = \mathrm{D}\) is a linear form on the left D-vector space \(\operatorname{Hom}_{\mathrm{A}}(\mathrm{S},\mathrm{M})\) .

PROPOSITION 6. --- Keep the notation above, and suppose that the A-module M is semisimple. The mapping \(u \mapsto h(u)\) from the right D-vector space \(\mathrm{Hom}_{\mathrm{A}}(\mathrm{M},\mathrm{S})\) to the dual of the left D-vector space \(\mathrm{Hom}_{\mathrm{A}}(\mathrm{S},\mathrm{M})\) is D-linear and bijective.

Let \(u \in \mathrm{Hom}_{\mathrm{A}}(\mathrm{M}, \mathrm{S})\) , \(v \in \mathrm{Hom}_{\mathrm{A}}(\mathrm{S}, \mathrm{M})\) , and \(d \in \mathrm{D}\) . We have

\[
h (u d) (v) = h (d \circ u) (v) = d \circ u \circ v = d \circ (h (u) (v)) = h (u) (v) d.
\]

This proves that the mapping h is D-linear. It is simply the mapping given by \(u \mapsto \operatorname{Hom}(1_{\mathrm{S}}, u)\) from \(\operatorname{Hom}_{\mathrm{A}}(\mathrm{M}, \mathrm{S})\) to \(\operatorname{Hom}_{\mathrm{D}}(\operatorname{Hom}_{\mathrm{A}}(\mathrm{S}, \mathrm{M}), \operatorname{Hom}_{\mathrm{A}}(\mathrm{S}, \mathrm{S}))\) . To prove that it is bijective, by Proposition 5, b) of VIII, p.~66, it suffices to treat the case when M is isotypical of type S; we can then apply the corollary of VIII, p.~65.

\subsection{Description of a Semisimple Module}

For the remainder of this section, A is a ring, and S is the set of classes of simple A-modules. For every \(\lambda \in S\) , we choose a simple module \(S_{\lambda}\) of class \(\lambda\) (for example, \(S_{\lambda} = \lambda\) ); we denote the opposite ring of the field of endomorphisms of \(S_{\lambda}\) by \(D_{\lambda}\) . We view \(S_{\lambda}\) as an \((A, D_{\lambda})\) -bimodule.

Let \(M\) be an A-module. For every \(\lambda \in \mathcal{S}\) , \(\mathrm{Hom}_{\mathrm{A}}(\mathrm{S}_{\lambda}, M)\) is a left vector space over the field \(D_{\lambda}\) . By VIII, p.~59, and II, §1, No.~6, p.~202, Proposition 6, there exists a unique A-linear mapping, called canonical,

\[
\alpha_ {\mathrm{M}}: \bigoplus_ {\lambda \in \mathcal {S}} \left(\mathrm{S} _ {\lambda} \otimes_ {\mathrm{D} _ {\lambda}} \operatorname{Hom} _ {\mathrm{A}} (\mathrm{S} _ {\lambda}, \mathrm{M})\right) \to \mathrm{M}
\]

satisfying the relation

\[
\alpha_ {\mathrm{M}} (s \otimes f) = f (s)\tag{10}
\]

for \(\lambda \in \mathcal{S}\) , \(s \in S_{\lambda}\) , and \(f \in \mathrm{Hom}_{\mathrm{A}}(\mathrm{S}_{\lambda}, \mathrm{M})\) . If we endow \(\bigoplus_{\lambda \in \mathcal{S}} (\mathrm{S}_{\lambda} \otimes_{\mathrm{D}_{\lambda}} \mathrm{Hom}_{\mathrm{A}}(\mathrm{S}_{\lambda}, \mathrm{M}))\) and \(\mathrm{M}\) with their natural structures of \(\mathrm{End}_{\mathrm{A}}(\mathrm{M})\) -modules, then the mapping \(\alpha_{\mathrm{M}}\) is \(\mathrm{End}_{\mathrm{A}}(\mathrm{M})\) -linear.

PROPOSITION 7. --- Let M be an A-module. The canonical mapping \(\alpha_{\mathrm{M}}\) is injective. For every \(\lambda \in \mathcal{S}\) , the image under \(\alpha_{\mathrm{M}}\) of \(\mathrm{S}_{\lambda} \otimes_{\mathrm{D}_{\lambda}} \mathrm{Hom}_{\mathrm{A}}(\mathrm{S}_{\lambda}, \mathrm{M})\) is the isotypical component of M of type \(\lambda\) . The image of \(\alpha_{\mathrm{M}}\) is the socle of M. The A-module M is semisimple if and only if the mapping \(\alpha_{\mathrm{M}}\) is bijective.

Let \(\lambda \in \mathcal{S}\) . Denote the isotypical component of M of type \(\lambda\) by \(\mathrm{M}_{\lambda}\) . Every A-linear mapping from \(\mathrm{S}_{\lambda}\) to M takes its values in \(\mathrm{M}_{\lambda}\) (VIII, p.~66, Proposition 5). Consequently, by Proposition 3, a) of VIII, p.~62, the mapping \(\alpha_{\mathrm{M}}\) induces a bijection from \(\mathrm{S}_{\lambda} \otimes_{\mathrm{D}_{\lambda}} \mathrm{Hom}_{\mathrm{A}}(\mathrm{S}_{\lambda}, \mathrm{M})\) to \(\mathrm{M}_{\lambda}\) . The proposition follows because the socle of M is the direct sum of the family \((\mathrm{M}_{\lambda})_{\lambda \in \mathcal{S}}\) and the module M is semisimple if and only if it is equal to its socle.

DEFINITION 5. --- Let M be a semisimple A-module. A description of M (with respect to the family \((\mathrm{S}_{\lambda})_{\lambda\in\mathcal{S}}\) ) is a pair \(((\mathrm{V}_{\lambda})_{\lambda\in\mathcal{S}},\alpha)\) , where \(V_{\lambda}\) is a left vector space over the field \(D_{\lambda}\) for each \(\lambda\in S\) and \(\alpha:\bigoplus_{\lambda\in\mathcal{S}}(\mathrm{S}_{\lambda}\otimes_{\mathrm{D}_{\lambda}}\mathrm{V}_{\lambda})\to\mathrm{M}\) is an isomorphism of A-modules.

By Proposition 7, every semisimple A-module M has a canonical description: it is the pair \(\left(\left(\mathrm{Hom}_{\mathrm{A}}\left(\mathrm{S}_{\lambda},\mathrm{M}\right)\right)_{\lambda \in \mathcal{S}},\alpha_{\mathrm{M}}\right)\) , where \(\alpha_{\mathrm{M}}\) is the A-linear mapping defined by formula (10).

PROPOSITION 8. --- Let M be a semisimple A-module and \(((\mathrm{V}_{\lambda})_{\lambda \in \mathcal{S}}, \alpha)\) a description of M.

\begin{enumerate}
\def\labelenumi{\alph{enumi})}
\item
  For every \(\lambda \in \mathcal{S}\) , the mapping \(\alpha\) induces an isomorphism from the A-module \(S_{\lambda} \otimes_{D_{\lambda}} V_{\lambda}\) to the isotypical component of M of type \(\lambda\) .
\item
  For every \(\lambda \in \mathcal{S}\) , the mapping \(\beta_{\lambda} : \mathrm{V}_{\lambda} \to \mathrm{Hom}_{\mathrm{A}}(\mathrm{S}_{\lambda}, \mathrm{M})\) defined by \(\beta_{\lambda}(v)(s) = \alpha(s \otimes v)\) is \(\mathrm{D}_{\lambda}\) -linear and bijective.
\item
  Let \(\mathbf{N}\) be a submodule of \(\mathbf{M}\) . There exists a unique family \((\mathrm{W}_{\lambda})_{\lambda \in \mathcal{S}}\) with the following properties: \(\mathrm{W}_{\lambda}\) is a \(\mathrm{D}_{\lambda}\) -linear subspace of \(\mathrm{V}_{\lambda}\) for every \(\lambda \in \mathcal{S}\) , and \(\mathbf{N}\) is the image under \(\alpha\) of the module \(\bigoplus_{\lambda \in \mathcal{S}} (\mathrm{S}_{\lambda} \otimes_{\mathrm{D}_{\lambda}} \mathrm{W}_{\lambda})\) identified with its canonical image in the module \(\bigoplus_{\lambda \in \mathcal{S}} (\mathrm{S}_{\lambda} \otimes_{\mathrm{D}_{\lambda}} \mathrm{V}_{\lambda})\) . For every \(\lambda \in \mathcal{S}\) , \(\mathrm{W}_{\lambda}\) is the set of elements \(v \in \mathrm{V}_{\lambda}\) such that \(\alpha(s \otimes v)\) belongs to \(\mathbf{N}\) for every \(s \in \mathrm{S}_{\lambda}\) .
\end{enumerate}

The A-module \(S_{\lambda} \otimes_{D_{\lambda}} V_{\lambda}\) is isotypical of type \(\lambda\) for every \(\lambda \in S\) . Assertion a) then follows from the facts that \(\alpha\) is an isomorphism and that M is the direct sum of the family \((M_{\lambda})_{\lambda \in \mathcal{S}}\) (VIII, p.~65, Proposition 4, a)).

Let \(\lambda \in \mathcal{S}\) . Denote the restriction of \(\alpha\) to \(S_{\lambda} \otimes_{D_{\lambda}} V_{\lambda}\) by \(\alpha_{\lambda}: S_{\lambda} \otimes_{D_{\lambda}} V_{\lambda} \to M\) . With the notation of No.~3 applied to the \((A, D_{\lambda})\) -bimodule \(S_{\lambda}\) , we have

\[
\beta_ {\lambda} = \gamma (\alpha_ {\lambda}) = \mathcal {H} (\alpha_ {\lambda}) \circ \beta_ {\mathrm{V} _ {\lambda}},
\]

where the last equality follows from VIII, p.~60. Hence, \(\beta_{\lambda}\) is the composition of the \(D_{\lambda}\) -linear homomorphism \(\mathcal{H}(\alpha_{\lambda})\) and \(\beta_{V_{\lambda}}\) . Assertion b) then follows from Proposition 3, b) of VIII, p.~62.

Let \(N\) be a submodule of \(M\) . We have \(N = \bigoplus_{\lambda \in \mathcal{S}} (N \cap M_{\lambda})\) (VIII, p.~66, Corollary), so c) follows from Theorem 2 of VIII, p.~62.

COROLLARY. --- Let M be a semisimple A-module. For every submodule N of M and every element \(\lambda\) of S, identify \(\operatorname{Hom}_{\mathrm{A}}(\mathrm{S}_{\lambda},\mathrm{N})\) with the \(D_{\lambda}\) -linear subspace of \(\operatorname{Hom}_{\mathrm{A}}(\mathrm{S}_{\lambda},\mathrm{M})\) consisting of the mappings with image contained in N.

\begin{enumerate}
\def\labelenumi{\alph{enumi})}
\item
  The mapping \(\mathrm{N}\mapsto(\mathrm{Hom}_{\mathrm{A}}(\mathrm{S}_{\lambda},\mathrm{N}))_{\lambda\in\mathcal{S}}\) is a bijection from the set of A-submodules of M to the set of families \((\mathrm{W}_{\lambda})_{\lambda\in\mathcal{S}}\) such that for every \(\lambda\in S\) , \(W_{\lambda}\) is a \(D_{\lambda}\) -linear subspace of \(\mathrm{Hom}_{\mathrm{A}}(\mathrm{S}_{\lambda},\mathrm{M})\) .
\item
  The inverse bijection sends a family \((\mathrm{W}_{\lambda})_{\lambda \in \mathcal{S}}\) to the A-submodule \(\sum_{\lambda \in \mathcal{S}}\sum_{w\in \mathrm{W}_{\lambda}}w(\mathrm{S}_{\lambda})\) of M.
\end{enumerate}

This is a reformulation of Proposition 8, c) applied to the canonical description of M.

PROPOSITION 9. --- Let M and M' be semisimple A-modules, and consider descriptions \(((\mathrm{V}_{\lambda})_{\lambda\in\mathcal{S}},\alpha)\) and \(((\mathrm{V}_{\lambda}^{\prime})_{\lambda\in\mathcal{S}},\alpha^{\prime})\) of M and M', respectively. For every family \(\boldsymbol{f}=(f_{\lambda})_{\lambda\in\mathcal{S}}\) in \(\prod_{\lambda\in\mathcal{S}}\mathrm{Hom}_{\mathrm{D}_{\lambda}}(\mathrm{V}_{\lambda},\mathrm{V}_{\lambda}^{\prime})\) , there exists a unique A-linear mapping \(\varphi(\boldsymbol{f})\in\mathrm{Hom}_{\mathrm{A}}(\mathrm{M},\mathrm{M}^{\prime})\) for which the following diagram commutes:

\[
\begin{array}{c} \bigoplus_ {\lambda \in \mathscr {S}} (S _ {\lambda} \otimes_ {D _ {\lambda}} V _ {\lambda}) \xrightarrow {\alpha} M \\ \bigoplus_ {\lambda \in \mathscr {S}} (1 _ {S _ {\lambda}} \otimes f _ {\lambda}) \Bigg | _ {\downarrow} \qquad \qquad \qquad \qquad \qquad \qquad \Bigg | _ {\varphi (\boldsymbol {f})} \\ \bigoplus_ {\lambda \in \mathscr {S}} (S _ {\lambda} \otimes_ {D _ {\lambda}} V _ {\lambda} ^ {\prime}) \xrightarrow {\alpha^ {\prime}} M ^ {\prime}. \end{array}
\]

The mapping \(\varphi : \prod_{\lambda \in \mathcal{S}} \mathrm{Hom}_{\mathrm{D}_{\lambda}}(\mathrm{V}_{\lambda}, \mathrm{V}_{\lambda}') \to \mathrm{Hom}_{\mathrm{A}}(\mathrm{M}, \mathrm{M}')\) defined this way is a group isomorphism. When we have \(\mathrm{M} = \mathrm{M}'\) , \(\mathrm{V}_{\lambda} = \mathrm{V}_{\lambda}'\) for every \(\lambda \in \mathcal{S}\) , and \(\alpha = \alpha'\) , the mapping \(\varphi\) is a ring isomorphism from \(\prod_{\lambda \in \mathcal{S}} \mathrm{End}_{\mathrm{D}_{\lambda}}(\mathrm{V}_{\lambda})\) to \(\mathrm{End}_{\mathrm{A}}(\mathrm{M})\) .

In view of the description of the isotypical components of M and M' given in Proposition 8, a), this follows from Theorem 3 of VIII, p.~64 and Proposition 5, b) of VIII, p.~66.

COROLLARY. --- Let M be a semisimple A-module and M' an A-module. The mapping \(u \mapsto (\operatorname{Hom}(1_{\mathrm{S}_{\lambda}}, u))_{\lambda \in \mathcal{S}}\) from \(\operatorname{Hom}_{\mathrm{A}}(\mathrm{M}, \mathrm{M}')\) to

\[
\prod_ {\lambda \in \mathcal {S}} \mathrm{Hom} _ {\mathrm{D} _ {\lambda}} (\mathrm{Hom} _ {\mathrm{A}} (\mathrm{S} _ {\lambda}, \mathrm{M}), \mathrm{Hom} _ {\mathrm{A}} (\mathrm{S} _ {\lambda}, \mathrm{M} ^ {\prime}))
\]

is a group isomorphism. When \(\mathrm{M}'\) is equal to \(\mathrm{M}\) , it is an isomorphism from the ring \(\operatorname{End}_{\mathrm{A}}(\mathrm{M})\) to the ring \(\prod_{\lambda \in \mathcal{S}} \operatorname{End}_{\mathrm{D}_{\lambda}}(\operatorname{Hom}_{\mathrm{A}}(\mathrm{S}_{\lambda}, \mathrm{M}))\) .

This is a reformulation of Proposition 9 applied to the canonical descriptions of M and of the socle of M'.

\subsection{Multiplicities and Lengths in Semisimple Modules}

PROPOSITION 10. --- Let M be a semisimple A-module. Let \((\mathrm{M}_{i})_{i\in\mathrm{I}}\) be a family of simple submodules with direct sum M. The following properties are equivalent:

\begin{enumerate}
\def\labelenumi{(\roman{enumi})}
\item
  M has finite length.
\item
  M is Artinian.
\item
  M is Noetherian.
\item
  M is finitely generated.
\item
  I is finite.
\end{enumerate}

If M has these properties, then the length of M is equal to the cardinal of I.

If the set I is finite, then M has properties (i), (ii), (iii), and (iv). Suppose that the set I is infinite. By Example 2 of VIII, p.~2, the module M is neither Artinian nor Noetherian; because every module of finite length is Artinian and Noetherian (VIII, p.~2, Proposition 1), M also does not have finite length. Finally, every element of M belongs to the sum of a finite number of submodules \(M_{i}\) , so M is not finitely generated. This proves the equivalence of properties (i) through (v). If these hold, then we have \(\text{long}(M) = \sum_{i \in I} \text{long}(M_{i}) = \text{Card}(I)\) (II, §1, No.~10, p.~213, Corollary 5).

PROPOSITION 11. --- Let M be a semisimple A-module that is the direct sum of a family \((\mathrm{M}_{i})_{i\in\mathrm{I}}\) of simple submodules. For every \(\lambda\in S\) , we denote by \(\mathrm{I}(\lambda)\) the set of indices \(i\in I\) such that \(M_{i}\) is of class \(\lambda\) . The cardinal of \(\mathrm{I}(\lambda)\) is equal to the dimension of the left \(D_{\lambda}\) -vector space \(\mathrm{Hom}_{\mathrm{A}}(\mathrm{S}_{\lambda},\mathrm{M})\) .

The isotypical component of A of type \(\lambda\) is isomorphic to \(\mathrm{S}_{\lambda}^{(\mathrm{I}(\lambda))}\) (VIII, p.~65, Proposition 4, b)). The \(D_{\lambda}\) -vector space \(\mathrm{Hom}_{\mathrm{A}}(\mathrm{S}_{\lambda},\mathrm{M})\) may be identified with \(\mathrm{Hom}_{\mathrm{A}}(\mathrm{S}_{\lambda},\mathrm{M}_{\lambda})\) , so is isomorphic to \(\mathrm{D}_{\lambda}^{(\mathrm{I}(\lambda))}\) (No.~2). This proves the proposition.

Every simple module is primordial (VIII, p.~45), so every semisimple module is semiprimordial. Let M be a semisimple A-module, and let \(\lambda \in \mathcal{S}\) . The multiplicity of \(\lambda\) in M is the primordial multiplicity \([\mathrm{M}:\lambda]\) of \(\lambda\) in M defined in VIII, p.~34. Proposition 11 translates to the equality

\[
[ \mathrm{M}: \lambda ] = \dim_ {\mathrm{D} _ {\lambda}} (\operatorname{Hom} _ {\mathrm{A}} (\mathrm{S} _ {\lambda}, \mathrm{M})).\tag{11}
\]

More generally, if \(((\mathrm{V}_{\lambda})_{\lambda \in \mathcal{S}},\alpha)\) is a description of M, then \([\mathrm{M}:\lambda ]\) is equal to \(\dim_{\mathrm{D}_{\lambda}}(\mathrm{V}_{\lambda})\) . By Proposition 6 of VIII, p.~68, we also have

\[
[ \mathrm{M}: \lambda ] = \dim_ {\mathrm{D} _ {\lambda}} (\operatorname{Hom} _ {\mathrm{A}} (\mathrm{M}, \mathrm{S} _ {\lambda}))\tag{12}
\]

when the multiplicity \([M : \lambda]\) is finite. Semisimple A-modules M and \(M'\) are isomorphic if and only if we have \([M : \lambda] = [M' : \lambda]\) for every \(\lambda \in S\) .

Let M be a semisimple A-module. There exists a cardinal I that has the following property: for every decomposition \(M = \bigoplus_{i \in I} M_i\) of M into a direct sum of simple modules, the cardinal of I is equal to I (VIII, p.~34, Corollary 2). This cardinal is called the length of the semisimple A-module M and denoted by \(\operatorname{long}_{\mathrm{A}}(\mathrm{M})\) or \(\operatorname{long}(\mathrm{M})\) . When M has finite length, this definition is compatible with that of II, §1, No.~10, p.~212, by Proposition 10.

The simple A-modules are the semisimple A-modules of length 1, and we have the formula

\[
\operatorname{long} _ {\mathrm{A}} \left(\bigoplus_ {j \in \mathrm{J}} \mathrm{M} _ {j}\right) = \sum_ {j \in \mathrm{J}} \operatorname{long} _ {\mathrm{A}} (\mathrm{M} _ {j})\tag{13}
\]

for every family \((\mathrm{M}_j)_{j\in \mathrm{J}}\) of semisimple A-modules. By Proposition 11, we have

\[
\mathrm{long} _ {\mathrm{A}} (\mathrm{M}) = \sum_ {\lambda \in \mathcal {S}} \dim_ {\mathrm{D} _ {\lambda}} \mathrm{Hom} _ {\mathrm{A}} (\mathrm{S} _ {\lambda}, \mathrm{M}).\tag{14}
\]

By applying this formula to \(M_{\lambda}\) , we obtain \([M : \lambda] = \text{long}_{A}(M_{\lambda})\) for every \(\lambda \in S\) .

When A is a field, the simple A-modules are the vector spaces of dimension 1; every A-module is then semisimple (II, §7, No.~1, p.~292, Theorem 1), and its length is simply its dimension as a vector space over A (II, §7, No.~2, p.~293).

\BourbakiExercises{Exercises}

\begin{enumerate}
\def\labelenumi{\arabic{enumi})}
\tightlist
\item
  \begin{enumerate}
  \def\labelenumii{\alph{enumii})}
  \tightlist
  \item
    Let M be an A-module and x an element of M whose annihilator is the intersection of a finite family of maximal ideals. Prove that the module Ax is semisimple (embed it in a finite product of simple modules).
  \end{enumerate}
\end{enumerate}

\begin{enumerate}
\def\labelenumi{\alph{enumi})}
\setcounter{enumi}{1}
\tightlist
\item
  Deduce that a module M is semisimple if and only if the annihilator of every nonzero element is the intersection of a finite family of maximal ideals.
\end{enumerate}

\begin{enumerate}
\def\labelenumi{\arabic{enumi})}
\setcounter{enumi}{1}
\tightlist
\item
  Let M be a module. Prove that the following properties are equivalent:
\end{enumerate}

\(\alpha)\) M is semisimple and has finite length.

\(\beta)\) M is Noetherian, and every maximal submodule of M admits a supplement.

\(\gamma)\) M is Artinian, and every simple submodule of M admits a supplement.

\begin{enumerate}
\def\labelenumi{\arabic{enumi})}
\setcounter{enumi}{2}
\item
  Let \(\mathrm{K}\) be a field, I an infinite set, and A the ring \(\mathrm{K}^{\mathrm{I}}\) . Prove that the A-module \(\mathrm{A}_s\) is not semisimple even though every monogenous ideal admits a supplement (determine all simple submodules of \(\mathrm{A}_s\) ).
\item
  Let M be a module that is the direct sum of a family \((\mathrm{M}_{\alpha})_{\alpha\in\mathrm{I}}\) of submodules. Prove that the submodules \(M_{\alpha}\) are stable under all automorphisms of M if and only if every homomorphism \(M_{\alpha}\to M_{\beta}\) with \(\alpha\neq\beta\) is zero.
\item
  Let M be a module of finite length, and let \(M = \oplus_{i \in I} M_i\) be a decomposition of M into a direct sum of indecomposable submodules. Let \(J \subset I\) be an equivalence class for the relation ``the modules \(M_i\) and \(M_j\) are isomorphic'' between i and j; we set \(M_J = \oplus_{i \in J} M_i\) . Give an example to show that \(M_J\) is not necessarily stable under every automorphism of M (use Exercise 4 with M taken to be a Z-module of finite length).
\item
  A module M is called semi-Artinian if every nonzero quotient of M contains a simple submodule. Every Artinian module and every vector space over a field is semi-Artinian.
\end{enumerate}

\begin{enumerate}
\def\labelenumi{\alph{enumi})}
\item
  Let \(0 \rightarrow M' \rightarrow M \rightarrow M'' \rightarrow 0\) be an exact sequence. Then M is semi-Artinian if and only if the same holds for \(M'\) and \(M''\) (if M is semi-Artinian, consider a maximal element in the set of submodules N of M such that \(N \cap M' = 0\) ).
\item
  The direct sum of a family of semi-Artinian modules is a semi-Artinian module. c) Let K be a field and A the product ring \(K^{N}\) . Prove that the A-module \(A_{s}\) is not semi-Artinian even though it is a product of simple A-modules (if s is the socle of \(A_{s}\) , prove that the A-module \(A_{s}/s\) does not contain any simple modules).
\item
  Let A be a left Noetherian ring and M a semi-Artinian A-module. Prove that M is the sum of its Artinian submodules (let S be the sum of the Artinian submodules of M; observe that if N is a finitely generated submodule of M with simple image in M/S, then \(N \cap S\) is Artinian). In particular, if M is finitely generated, then it is Artinian.
\item
  Let D be a field, V an infinite-dimensional D-vector space, and A the subring of \(\operatorname{End}_{\mathrm{D}}(\mathrm{V})\) generated by the center of D and the endomorphisms of finite rank. Prove that \(A_{s}\) is semi-Artinian (use a) and b)) but is not the sum of its Artinian submodules.
\end{enumerate}

\begin{enumerate}
\def\labelenumi{\arabic{enumi})}
\setcounter{enumi}{6}
\tightlist
\item
  Let A be a ring and M an A-module. A submodule N of M is called essential if \(N \cap P \neq 0\) for every nonzero submodule P of M.
\end{enumerate}

\begin{enumerate}
\def\labelenumi{\alph{enumi})}
\item
  Prove that every submodule N of M is a direct factor of an essential submodule (consider a submodule P that is maximal among the submodules of A with zero intersection with N).
\item
  Prove that the socle of M is the intersection of the essential submodules of M (deduce from a) that this intersection is semisimple). In particular, M is semisimple if and only if it does not contain any essential submodules other than itself.
\end{enumerate}

¶ 8) Let A be a left Noetherian ring.

\begin{enumerate}
\def\labelenumi{\alph{enumi})}
\item
  Let x be a nonzero element of A. Prove that the left annihilator Ann x of x is not essential (among the elements that do not have this property, take an x such that Ann x is maximal, and consider \(Ax \cap Ann x\) ).
\item
  Let x be a right cancellable element of A. Prove that the ideal Ax is essential (if a is an ideal of A such that \(a \cap Ax = 0\) , consider the ideal \(a + ax + ax^{2} + \ldots\) ).
\item
  Deduce from a) and b) that every right cancellable element of A is cancellable.
\item
  Let a be an essential left ideal of A containing a nonnilpotent element. Prove that a contains a cancellable element (let \((x_{1},\ldots,x_{n})\) be a sequence of nonzero elements of a satisfying Ann \(x_{i}^{2}=\) Ann \(x_{i}\) and \(x_{i+1}\in\) Ann \(x_{i}\cap\cdots\cap\) Ann \(x_{1}\) for every \(i\geqslant1\) ; observe that a contains the direct sum of the \(Ax_{i}\) . Deduce that there exists such a sequence with \(\bigcap_{i}\) Ann \(x_{i}=0\) ; prove that the element \(\sum x_{i}^{2}\) is then right cancellable, and use c)).
\end{enumerate}

\section{Commutation}

\subsection{The Commutant and Bicommutant of a Module}

Let E be a ring and B a subset of E. The commutant (or centralizer) of B in E is the subring \(B'\) of E consisting of the elements that commute with every element of B. The commutant \(B''\) of \(B'\) is called the bicommutant (or bicentralizer) of B. We have \(B \subset B''\) , and \(B'\) coincides with its bicommutant (III, §1, No.~2, p.~429). The center of a subring B of E is \(B \cap B'\) ; the common center of \(B'\) and \(B''\) is \(B' \cap B''\) . If B is a commutative subring of E, then we have \(B \subset B'\) , and \(B''\) is the center of \(B'\) (III, §1, No.~2, p.~430).

Let A be a ring and M a left (resp. right) A-module. Let us apply these definitions to the case when E is the endomorphism ring of the additive group of M and B the ring of homotheties \(A_{M}\) of M. The commutant \(A_{M}^{\prime}\) of \(A_{M}\) in E is called the commutant of M; it is the endomorphism ring of the A-module M. The bicommutant \(A_{M}^{\prime\prime}\) of \(A_{M}\) in E is called the bicommutant of M; it is the endomorphism ring of the countermodule of M (VIII, p.~8, Definition 3). We have \(A_{M} \subset A_{M}^{\prime\prime}\) , the ring \(A_{M} \cap A_{M}^{\prime}\) is the center of \(A_{M}\) , and \(A_{M}^{\prime} \cap A_{M}^{\prime\prime}\) is the common center of \(A_{M}^{\prime}\) and \(A_{M}^{\prime\prime}\) .

DEFINITION 1. --- The A-module M is balanced if we have \(\mathrm{A}_{\mathrm{M}} = \mathrm{A}_{\mathrm{M}}^{\prime \prime}\) .

If the A-module M is balanced, then the rings \(A_{M}\) and \(A_{M}^{\prime}\) have the same center \(A_{M} \cap A_{M}^{\prime}\) . The module M is a balanced A-module if and only if it is a balanced \(A_{M}\) -module. The countermodule of an A-module M is faithful and balanced.

When the ring A is commutative, the bicommutant \(A_{M}^{\prime\prime}\) of M is the center of \(\mathrm{A}_{\mathrm{M}}^{\prime}=\mathrm{End}_{\mathrm{A}}(\mathrm{M})\) ; that M is balanced means that the center of \(\mathrm{End}_{\mathrm{A}}(\mathrm{M})\) is reduced to the homotheties.

For any element \(a\) of A, denote by \(\delta_{a}\) the right homothety \(x\mapsto xa\) from A to A and by \(\gamma_{a}\) the left homothety \(x\mapsto ax\) (I, §8, No.~1, p.~97). The map \(a\mapsto \delta_{a}\) is a ring isomorphism from \(\mathrm{A}^{\circ}\) to the commutant of the left A-module \(\mathrm{A}_s\) (cf.~II, §10, No.~7, p.~349, applied to \(\mathrm{I} = \{1\}\) ). The map \(a\mapsto \gamma_{a}\) is a ring isomorphism from A to the commutant of the right A-module \(\mathrm{A}_d\) (loc. cit.). If we identify A with the commutant of \(\mathrm{A}_d\) through this map, then the countermodule of \(\mathrm{A}_d\) is identified with \(\mathrm{A}_s\) ; consequently, the A-module \(\mathrm{A}_d\) is balanced. Likewise, the A-module \(\mathrm{A}_s\) is balanced.

Let n be an integer \(\geqslant 1\) . We view \(A^{n}\) as a left \(\mathbf{M}_{n}(\mathrm{A})\) -module (loc. cit.). The mapping that sends \(m \in \mathbf{M}_{n}(\mathrm{A})\) to the endomorphism \(x \mapsto mx\) of the A-module \(A_{d}^{n}\) is a ring isomorphism from \(\mathbf{M}_{n}(\mathrm{A})\) to the commutant of the right A-module \(A_{d}^{n}\) (loc. cit.).

PROPOSITION 1. --- Let \((\mathrm{A}_{i})_{i\in\mathrm{I}}\) be a family of rings, and for every \(i\in I\) , let \(M_{i}\) be an \(A_{i}\) -module. Set \(A=\prod A_{i}, M=\prod M_{i}\) , and \(N=\bigoplus M_{i}\) . Endow M with the structure of an A-module with law of action \(((a_{i}),(x_{i}))\mapsto(a_{i}x_{i})\) . The set N is an A-submodule of M.

\begin{enumerate}
\def\labelenumi{\alph{enumi})}
\item
  The mapping \((u_i) \mapsto \prod u_i\) from \(\prod \mathrm{End}_{\mathbf{Z}}(\mathrm{M}_i)\) to \(\mathrm{End}_{\mathbf{Z}}(\mathrm{M})\) (II, §1, No.~5, p.~200) restricts to ring isomorphisms from \(\prod (\mathrm{A}_i)_{\mathrm{M}_i}\) , \(\prod (\mathrm{A}_i)'_{\mathrm{M}_i}\) , and \(\prod (\mathrm{A}_i)'''_{\mathrm{M}_i}\) to \(\mathrm{A}_{\mathrm{M}}\) , \(\mathrm{A}_{\mathrm{M}}'\) , and \(\mathrm{A}_{\mathrm{M}}''\) , respectively.
\item
  The mapping \((u_i) \mapsto \bigoplus u_i\) from \(\prod \operatorname{End}_{\mathbf{Z}}(\mathrm{M}_i)\) to \(\operatorname{End}_{\mathbf{Z}}(\mathrm{N})\) (II, §1, No.~6, p.~203) restricts to ring isomorphisms from \(\prod (\mathrm{A}_i)_{\mathrm{M}_i}\) , \(\prod (\mathrm{A}_i)'_{\mathrm{M}_i}\) , and \(\prod (\mathrm{A}_i)'''_{\mathrm{M}_i}\) to \(\mathrm{A_N}\) , \(\mathrm{A_N}'\) , and \(\mathrm{A_N}''\) , respectively.
\end{enumerate}

The ring isomorphisms defined in Proposition 1 are called canonical. The product \(\prod (\mathrm{A}_i)_{\mathrm{M}_i}\) is often identified with \(\mathrm{A_M}\) , and \(\prod (\mathrm{A}_i)'_{\mathrm{M}_i}\) with \(\mathrm{A_M}'\) , etc., through these isomorphisms.

The mapping \(\varphi : (u_i) \mapsto \prod u_i\) from \(\prod \operatorname{End}_{\mathbf{Z}}(\mathrm{M}_i)\) to \(\operatorname{End}_{\mathbf{Z}}(\mathrm{M})\) is an injective ring homomorphism. By the definition of the A-module structure on M, we have \(\varphi(\prod (\mathrm{A}_i)_{\mathrm{M}_i}) = \mathrm{A}_{\mathrm{M}}\) . Let \(u \in \mathrm{A}_{\mathrm{M}}'\) . For every \(i \in \mathrm{I}\) , denote by \(h_i\) the element of A with all components equal to 1 except that of index \(i\) , which is equal to 0. If \(x\) is an element of M with component of index \(i\) equal to 0, then we have \(x = h_i x\) , and therefore \(\operatorname{pr}_i(u(x)) = \operatorname{pr}_i(u(h_i x)) = \operatorname{pr}_i(h_i u(x)) = 0\) . Consequently, there exists a unique group homomorphism \(u_i: \mathrm{M}_i \mapsto \mathrm{M}_i\) such that \(\operatorname{pr}_i(u(y)) = u_i(\operatorname{pr}_i(y))\) for every \(y \in \mathrm{M}\) . We have \(u = \prod u_i\) . Since the mapping \(u\) is A-linear, the mapping \(u_i\) is \(\mathrm{A}_i\) -linear for every \(i \in \mathrm{I}\) . This proves that the image of \(\prod (\mathrm{A}_i)'_{\mathrm{M}_i}\) under \(\varphi\) contains \(\mathrm{A}_{\mathrm{M}}'\) ; the reverse inclusion is obvious. By applying this to the countermodule of M, we deduce that \(\varphi\) induces an isomorphism from \(A_{M}^{\prime \prime}\) to \(\prod_i(A_i)_{M_i}^{\prime \prime}\) . This proves assertion a).

The proof of b) is the same as that of a) mutatis mutandis.

PROPOSITION 2. --- Let A be a ring, and let M be an A-module. Let I be a set. Then the bicommutant of the A-module \(\mathrm{M}^{(I)}\) coincides with the ring of homotheties of the \(A_{M}^{\prime\prime}\) -module \(\mathrm{M}^{(I)}\) .

For \(i \in I\) , we denote the projection homomorphism from \(M^{(I)}\) to M by \(\pi_{i} : (x_{j})_{j \in I} \mapsto x_{i}\) and the corresponding canonical injection (I, §4, No.~9, p.~47) by \(\iota_{i} : M \to M^{(I)}\) .

Let \(u\) be an element of \(\mathrm{End}_{\mathrm{A}}(\mathrm{M}^{(\mathrm{I})})\) . For all \(i,j \in \mathrm{I}\) , the composition \(u_{i,j} = \pi_j \circ u \circ \iota_i\) belongs to the commutant \(\mathrm{A}_{\mathrm{M}}'\) of \(\mathrm{M}\) . For every element \(b\) of \(\mathrm{A}_{\mathrm{M}}''\) and every \((x_i) \in \mathrm{M}^{(\mathrm{I})}\) , we have the relations

\[
b u \left(\left(x _ {i}\right) _ {i \in \mathrm{I}}\right) = b \left(\sum_ {i \in \mathrm{I}} u _ {i, j} \left(x _ {i}\right)\right) _ {j \in \mathrm{I}} = \left(\sum_ {i \in \mathrm{I}} u _ {i, j} \left(b x _ {i}\right)\right) _ {j \in \mathrm{I}} = u \left(b \left(x _ {i}\right) _ {i \in \mathrm{I}}\right).
\]

The homothety \(b_{\mathrm{M}^{(\mathrm{I})}}\) therefore belongs to the bicommutant of the A-module \(\mathrm{M}^{(\mathrm{I})}\) .

Conversely, let b be an element of the bicommutant of \(\mathrm{M}^{(\mathrm{I})}\) . For \(i,j\in I\) , set \(b_{i,j}=\pi_{j}\circ b\circ\iota_{i}\) . Take \(i,j\in I\) with \(i\neq j\) . Since \(\iota_{j}\circ\pi_{j}\) belongs to the commutant of the A-module \(\mathrm{M}^{(\mathrm{I})}\) , we have

\[
b _ {i, j} = \pi_ {j} \circ \iota_ {j} \circ \pi_ {j} \circ b \circ \iota_ {i} = \pi_ {j} \circ b \circ \iota_ {j} \circ \pi_ {j} \circ \iota_ {i} = 0.
\]

Likewise, we have

\[
b _ {j, j} = \pi_ {j} \circ b \circ \iota_ {j} = \pi_ {i} \circ \iota_ {i} \circ \pi_ {j} \circ b \circ \iota_ {j} = \pi_ {i} \circ b \circ \iota_ {i} \circ \pi_ {j} \circ \iota_ {j} = b _ {i, i}.
\]

Moreover, \(b_{i,i}\) belongs to \(A_{M}^{\prime\prime}\) . It follows that b coincides with a homothety of the \(A_{M}^{\prime\prime}\) -module \(M^{(I)}\) .

\subsection{Generating Modules}

Let \(\mathrm{A}\) be a ring.

DEFINITION 2. --- An A-module M is called generating if every A-module N is generated by the images of the A-linear mappings from M to N.

Let \(M\) be a left A-module. We denote the dual of \(M\) by \(M^*\) and the canonical bilinear form on \(M \times M^*\) (II, §2, No.~3, p.~233) by

\[
(x, x ^ {*}) \mapsto \langle x, x ^ {*} \rangle = x ^ {*} (x).
\]

We denote by \(\tau(\mathrm{M})\) the set of elements of A of the form \(\sum_{i=1}^{n}\langle x_i,x_i^*\rangle\) , where \(x_1,\ldots ,x_n\) are elements of M and \(x_1^*,\ldots ,x_n^*\) are elements of \(\mathrm{M}^*\) . It is a two-sided ideal of A, called the trace ideal of M. The trace ideal of the A-module \(\mathrm{A}_s\) is A. The trace ideal of a module that is the direct sum of a family \((\mathrm{M}_i)_{i\in \mathrm{I}}\) of A-modules is the ideal \(\sum_{i\in \mathrm{I}}\tau(\mathrm{M}_i)\) . If M is a projective A-module, then it follows from Proposition 12 of II, §2, No.~6, p.~238, that we have \(\mathrm{M} = \tau(\mathrm{M})\mathrm{M}\) .

THEOREM 1. --- Let M be a left A-module. The following properties are equivalent:

\begin{enumerate}
\def\labelenumi{(\roman{enumi})}
\item
  The A-module M is generating.
\item
  For every A-module N, there exist a set I and a surjective A-linear mapping from \(M^{(I)}\) to N.
\item
  There exist an integer \(n \geqslant 0\) and a surjective A-linear mapping from \(\mathbf{M}^n\) to \(\mathrm{A}_s\) .
\item
  There exists an integer \(n \geqslant 0\) such that \(A_s\) is isomorphic to a direct factor submodule of \(M^n\) .
\item
  The trace ideal \(\tau (\mathbf{M})\) is equal to A.
\item
  There exist an integer \(n \geqslant 0\) , elements \(x_{1}, \ldots, x_{n}\) of M, and elements \(x_{1}^{*}, \ldots, x_{n}^{*}\) of M* satisfying \(\sum_{i=1}^{n} \langle x_{i}, x_{i}^{*} \rangle = 1\) .
\item
  \(\Rightarrow\) (ii): Suppose that M is generating, and let N be a left A-module. There exists a family \((u_{i})_{i\in\mathrm{I}}\) of A-linear mappings from M to N such that we have \(N=\sum_{i\in\mathrm{I}}u_{i}(\mathrm{M})\) . The mapping \((x_{i})\mapsto\sum_{i\in\mathrm{I}}u_{i}(x_{i})\) from \(\mathrm{M}^{(\mathrm{I})}\) to N is A-linear and surjective.
\item
  \(\Rightarrow\) (iii): Assume that property (ii) holds, and take \(N = A_s\) . We then have a set I and a surjective linear mapping \(u: M^{(I)} \to A_s\) . Since \(A_s\) is generated by the element 1, there exists a finite subset J of I such that \(u(M^{(J)}) = A_s\) , whence (iii).
\item
  \(\Rightarrow\) (iv): This follows from Proposition 21 of II, §1, No.~12, p.~218.
\item
  \(\Rightarrow\) (v): Let \(n \geqslant 1\) be an integer such that \(\mathrm{A}_s\) is isomorphic to a direct factor submodule of M. We have \(\mathrm{A} = \tau(\mathrm{A}_s) \subset \tau(\mathrm{M}^n) = \tau(\mathrm{M})\) and therefore \(\tau(\mathrm{M}) = \mathrm{A}\) .
\item
  \(\Rightarrow\) (vi): This is clear.
\item
  \(\Rightarrow\) (i): Let \(n\) be an integer \(\geqslant 0\) , \(x_{1},\ldots ,x_{n}\) elements of M, and \(x_{1}^{*},\ldots ,x_{n}^{*}\) elements of M \(^*\) satisfying \(\sum_{i = 1}^{n}\langle x_i,x_i^*\rangle = 1\) . Let N be a left A-module and \(y\) an element of N. The mappings \(u_{i}:x\mapsto \langle x,x_{i}^{*}\rangle y\) from M to N are A-linear, and we have \(y = \sum_{i = 1}^{n}u_{i}(x_{i})\) . This proves that M is a generating A-module.
\end{enumerate}

COROLLARY. --- A generating A-module is faithful.

Let \(a \in A\) satisfy \(aM = 0\) . Using the implication (i) \(\Rightarrow\) (iv) of Theorem 1, we obtain \(aA_s = 0\) , and therefore \(a = 0\) . The corollary follows (II, §1, No.~12, p.~219).

Examples. --- 1) The A-module \(A_s\) is generating.

\begin{enumerate}
\def\labelenumi{\arabic{enumi})}
\setcounter{enumi}{1}
\item
  Every nonzero free A-module is generating. More generally, every module with a generating quotient is itself generating.
\item
  Let M be a semisimple A-module whose countermodule is finitely generated. Then M is a generating \(A_{M}\) -module. Indeed, by Lemma 4 of VIII, p.~8, there exists a natural number m such that \((\mathrm{A}_{\mathrm{M}})_{s}\) is isomorphic to a submodule of \(M^{m}\) . Since \(M^{m}\) is a semisimple \(A_{M}\) -module, \((\mathrm{A}_{\mathrm{M}})_{s}\) is isomorphic to a direct factor submodule of \(M^{m}\) and M is a generating \(A_{M}\) -module (Theorem 1).
\item
  Let A be a principal ideal domain, and let P be a finitely generated A-module. There exist an integer \(n \geqslant 1\) and an increasing sequence of ideals \((\mathfrak{a}_{i})_{1 \leqslant i \leqslant n}\) of A such that P is isomorphic to the direct sum of the \(A/\mathfrak{a}_{i}\) (VII, §4, No.~4, p.~19, Theorem 2); the annihilator a of P is equal to \(a_{1}\) . Then P is a generating module over the ring A/a. If P is not a torsion module, then we have a = 0 and P is a generating A-module.
\end{enumerate}

Lemma 1. --- Let A be a commutative ring, M a finitely generated A-module, and Ann(M) its annihilator. Let \(\mathfrak{a}\) be an ideal of A. The following properties are equivalent:

\begin{enumerate}
\def\labelenumi{(\roman{enumi})}
\item
  aM = M.
\item
  Ann(M) + a = A.
\item
  There exists an element \(a\) of \(\mathfrak{a}\) such that \(am = m\) for every \(m \in \mathrm{M}\) . (i) \(\Rightarrow\) (ii): Let \((x_1, \ldots, x_n)\) be a generating family of the A-module M. If \(\mathfrak{aM} = \mathrm{M}\) , then each \(x_i\) can be written as \(\sum_{j=1}^{n} c_{ij} x_j\) , where the \(c_{ij}\) belong to \(\mathfrak{a}\) . Denote the matrix \((c_{ij})\) by \(C\) and the column matrix with entries \(x_1, \ldots, x_n\) by \(X\) . We have \((I_n - C)X = 0\) . Let \(d\) be the determinant and \(V\) the cofactor matrix of the matrix \(I_n - C\) . By formula (26) of III, §8, No.~6, p.~532, we have \(dX = {}^t V(I_n - C)X = 0\) , so that \(d \in \operatorname{Ann}(\mathrm{M})\) . On the other hand, since the \(c_{ij}\) belong to \(\mathfrak{a}\) , we have \(d \equiv 1 (\bmod \mathfrak{a})\) (III, §8, No.~5, p.~529, (18)), and therefore \(1 \in \operatorname{Ann}(\mathrm{M}) + \mathfrak{a}\) .
\item
  \(\Rightarrow\) (iii): Assuming that (ii) holds, there exist \(a\in \mathfrak{a}\) and \(b\in \mathrm{Ann}(\mathrm{M})\) such that \(a + b = 1\) . We then have \(am = m\) for every \(m\in \mathrm{M}\) .
\item
  \(\Rightarrow\) (i): This is clear.
\end{enumerate}

PROPOSITION 3. --- Let A be a commutative ring. Every finitely generated and faithful projective A-module is generating. More generally, a finitely generated projective A-module P is a generating \(A_{P}\) -module.

Let P be a finitely generated projective A-module. We have \(\tau(\mathrm{P})\mathrm{P}=\mathrm{P}\) (VIII, p.~80).

If the A-module P is faithful, then the ideal \(\tau(\mathrm{P})\) is equal to A (Lemma 1) and the A-module P is generating (Theorem 1); this leads to the first assertion. The second follows by Lemma 2 below.

Lemma 2. --- Let A be a ring and M a projective A-module. The \(\mathrm{A_M}\) -module M is projective.

Let \((x_{i})_{i\in\mathrm{I}}\) be a generating family of the A-module M. There exists a family \((x_{i}^{*})_{i\in\mathrm{I}}\) of linear forms on the A-module M such that for every \(x\in\mathrm{M}\) , the family \((\langle x,x_{i}^{*}\rangle)_{i\in\mathrm{I}}\) has finite support and \(x=\sum_{i\in\mathrm{I}}\langle x,x_{i}^{*}\rangle x_{i}\) (II, §2, No.~6, p.~238, Proposition 12). For every \(i\in\mathrm{I}\) , the mapping \(x\mapsto\langle x,x_{i}^{*}\rangle_{\mathrm{M}}\) is a linear form on the \(A_{M}\) -module M, and we have \(x=\sum_{i\in\mathrm{I}}\langle x,x_{i}^{*}\rangle_{\mathrm{M}}x_{i}\) for every \(x\in\mathrm{M}\) . By loc. cit., M is a projective \(A_{M}\) -module.

\subsection{The Bicommutant of a Generating Module}

THEOREM 2. --- A generating module is balanced.

Let A be a ring, and let M be a generating A-module; by definition, there exist an integer \(n \geqslant 0\) , elements \(x_{1}, \ldots, x_{n}\) of M, and elements \(x_{1}^{*}, \ldots, x_{n}^{*}\) of the dual \(M^{*}\) of M satisfying \(\sum_{i=1}^{n} \langle x_{i}, x_{i}^{*} \rangle = 1\) . Recall (II, §4, No.~2, p.~271) that we define a group homomorphism \(\theta : M^{*} \otimes_{A} M \to \operatorname{End}_{A}(M)\) by the formula \(\theta(x^{*} \otimes y)(x) = \langle x, x^{*} \rangle y\) . If u is an element of the bicommutant of M, then it commutes with \(\operatorname{End}_{A}(M)\) ; therefore, for every \(y \in M\) , we have

\[
\begin{array}{l} u (y) = u \Big (\sum_ {i = 1} ^ {n} \langle x _ {i}, x _ {i} ^ {*} \rangle y \Big) = \sum_ {i = 1} ^ {n} u (\theta (x _ {i} ^ {*} \otimes y) (x _ {i})) \\ = \sum_ {i = 1} ^ {n} \theta (x _ {i} ^ {*} \otimes y) (u (x _ {i})) = \Big (\sum_ {i = 1} ^ {n} \langle u (x _ {i}), x _ {i} ^ {*} \rangle \Big) y. \end{array}
\]

Consequently, u belongs to \(A_{M}\) , and M is balanced.

COROLLARY 1. --- A free module is balanced.

This is clear if the module is zero, and a nonzero free module is generating (VIII, p.~81, Example 2).

COROLLARY 2. --- Let A be a ring, and let n be an integer \(\geqslant 0\) . The center of \(\mathbf{M}_{n}(\mathrm{A})\) consists of the scalar matrices with entries in the center of A. We view \(A^{n}\) as a left \(\mathbf{M}_{n}(\mathrm{A})\) -module (II, §10, No.~7, p.~349). The endomorphisms of this module are the mappings \(x \mapsto xa\) , where a runs through A.

Let M be the right A-module \(A_{d}^{n}\) . It is balanced by Corollary 1. Consequently, the centers of \(A_{M}\) , \(A_{M}^{\prime}\) , and \(A_{M}^{\prime\prime}\) coincide. Corollary 2 then follows from the fact that \(A_{M}\) is identified with A and \(A_{M}^{\prime}\) with \(\mathbf{M}_{n}(\mathbf{A})\) .

Remark. --- Let M be an A-module. If the \(A_{M}\) -module M is generating, then we have \(A_{M} = A_{M}^{\prime\prime}\) by Theorem 2 applied to the \(A_{M}\) -module M, so that M is a balanced A-module.

COROLLARY 3. --- Every finitely generated projective module over a commutative ring is balanced.

Indeed, a finitely generated projective module M is a generating \(A_{M}\) -module by Proposition 3 of VIII, p.~82. The corollary therefore follows from the remark above.

COROLLARY 4. --- Every finitely generated module over a principal ideal domain is balanced.

Indeed, a finitely generated module M over a principal ideal domain A is a generating \(A_{M}\) -module (VIII, p.~81, Example 4).

COROLLARY 5. --- Let K be a commutative field, V a finite-dimensional vector space over K, and u and v endomorphisms of V. The following properties are equivalent:

\begin{enumerate}
\def\labelenumi{(\roman{enumi})}
\item
  There exists a polynomial \(\mathrm{P}\) in \(\mathrm{K}[\mathrm{X}]\) such that \(v = \mathrm{P}(u)\) .
\item
  The endomorphism v commutes with every endomorphism of V that commutes with u.
\end{enumerate}

Take A to be the ring K{[}X{]} and M to be the K{[}X{]}-module deduced from V and u (VII, §5, No.~1, p.~28). Assertion (i) means that \(v \in K[X]_{M}\) and (ii) that \(v \in K[X]_{M}'\) . Corollary 5 is therefore a particular case of Corollary 4.

PROPOSITION 4. --- A semisimple module with finitely generated countermodule is balanced.

This follows from Example 3 of VIII, p.~81 and the remark.

COROLLARY 1. --- Let \((\mathrm{S}_{i})_{i\in\mathrm{I}}\) be a finite family of pairwise nonisomorphic simple A-modules. For \(i\in I\) , denote by \(D_{i}\) the opposite ring of the field of endomorphisms of \(S_{i}\) . Suppose that for every \(i\in I\) , the \(D_{i}\) -vector space \(S_{i}\) is finite-dimensional. Then the mapping \(a \mapsto (a_{\mathrm{S}_i})_{i \in \mathrm{I}}\) from A to \(\prod_{i \in \mathrm{I}} \mathrm{End}_{\mathrm{D}_i}(\mathrm{S}_i)\) is surjective.

Consider the A-module \(M = \prod_{i \in I} S_i\) . Since I is finite, we also have \(M = \bigoplus_{i \in I} S_i\) , and the image of \(S_i\) in M is the isotypical component of M of type \(S_i\) . Consequently, the endomorphisms of the A-module M are the mappings \((s_i) \mapsto (s_i d_i)\) , where \((d_i)_{i \in I}\) runs through \(\prod_{i \in I} D_i\) (VIII, p.~66, Proposition 5). Since I is finite and for every \(i \in I\) , the right \(D_i\) -vector space \(S_i\) is finite-dimensional, the countermodule of M is finitely generated. By Proposition 4, the A-module M is balanced. Now, the bicommutant of the A-module M consists of elements of \(\text{End}_{\mathbf{Z}}(\text{M})\) of the form \(\prod_{i \in I} u_i\) , where \((u_i) \in \prod_{i \in I} \text{End}_{\text{D}_i}(\text{S}_i)\) (VIII, p.~78, Proposition 1) because \(\text{End}_{\text{D}_i}(\text{S}_i)\) is the bicommutant of \(S_i\) . The corollary follows from this.

This corollary applies, in particular, when A is an algebra over a commutative field K and each \(S_{i}\) is a simple A-module that is finite-dimensional as a K-vector space: indeed, \(D_{i}\) then contains the homotheties \(\alpha_{S_{i}}\) , where \(\alpha\) runs through K, and \(S_{i}\) is finite-dimensional over \(D_{i}\) because it is so over K.

COROLLARY 2 (Burnside's theorem). --- Let A be an algebra over an algebraically closed commutative field K, and let S be a simple A-module that is finite-dimensional as a K-vector space. Then the mapping \(a \mapsto a_{\mathrm{S}}\) from A to \(\operatorname{End}_{\mathrm{K}}(\mathrm{S})\) is surjective.

Indeed, the field of endomorphisms of the A-module S consists of the homotheties \(\alpha_{S}\) with \(\alpha \in K\) (VIII, p.~47, Theorem 1). We then apply Corollary 1 to the simple A-module S.

\subsection{The Countermodule of a Semisimple Module}

Let A be a ring. Denote the set of classes of simple A-modules by S. For every \(\lambda \in S\) , choose a simple A-module \(S_{\lambda}\) of class \(\lambda\) , and denote the opposite ring of the field of endomorphisms of \(S_{\lambda}\) by \(D_{\lambda}\) . We view \(S_{\lambda}\) as an \((A, D_{\lambda})\) -bimodule.

Let \(M\) be a semisimple A-module, and let \(B\) be the endomorphism ring of \(M\) . Denote the bicommutant of \(M\) by \(C\) . For every \(\lambda \in \mathcal{S}\) , denote the left \((D_{\lambda}, B)\) -bimodule \(\mathrm{Hom}_A(S_{\lambda}, M)\) by \(V_{\lambda}\) . Finally, denote the support of the A-module \(M\) by \(\mathcal{S}_M\) (VIII, p.~66); it is also the set of elements \(\lambda\) of \(\mathcal{S}\) such that \(V_{\lambda}\) is nonzero.

Remark 1. --- The canonical description \(\alpha_{M}\) of the A-module M is an isomorphism of left (A, B)-bimodules. By the corollary of Proposition 9, VIII, p.~71, the mapping \(f \mapsto (\operatorname{Hom}(1_{\mathrm{S}_{\lambda}}, f))_{\lambda \in \mathcal{S}_{\mathrm{M}}}\) from B to \(\prod_{\lambda \in \mathcal{S}_{\mathrm{M}}} \operatorname{End}_{\mathrm{D}_{\lambda}}(\mathrm{V}_{\lambda})\) is a ring isomorphism.

PROPOSITION 5. --- a) The countermodule of M is semisimple.

\begin{enumerate}
\def\labelenumi{\alph{enumi})}
\setcounter{enumi}{1}
\item
  For every \(\lambda \in \mathcal{S}_{\mathrm{M}}\) , the B-module \(V_{\lambda}\) is simple, and its commutant is equal to \((D_{\lambda})_{V_{\lambda}}\) .
\item
  The mapping \(\lambda \mapsto \mathrm{cl}(\mathrm{V}_{\lambda})\) is a bijection from the support of the A-module M to the support of its countermodule.
\item
  For every \(\lambda \in \mathcal{S}_{\mathrm{M}}\) , the B-submodule \(\mathrm{M}_{\lambda}\) is the isotypical component of type \(\mathrm{V}_{\lambda}\) of the B-module M, and the multiplicity of \(\mathrm{V}_{\lambda}\) in M is equal to \(\dim_{\mathrm{D}_{\lambda}}(\mathrm{S}_{\lambda})\) .
\item
  For \(s \in S\) , denote the mapping \(\varphi \mapsto \varphi(s)\) from \(V_{\lambda} = \operatorname{Hom}_{\mathrm{A}}(\mathrm{S}_{\lambda}, \mathrm{M})\) to M by \(\widetilde{s}\) . It is B-linear. The resulting mapping \(s \mapsto \widetilde{s}\) from \(S_{\lambda}\) to \(\operatorname{Hom}_{\mathrm{B}}(\mathrm{V}_{\lambda}, \mathrm{M})\) is an isomorphism of \((\mathrm{A}, \mathrm{D}_{\lambda})\) -bimodules.
\end{enumerate}

Let \(\lambda \in \mathcal{S}_{\mathrm{M}}\) . Denote the ring \(\mathrm{End}_{\mathrm{D}_{\lambda}}(\mathrm{V}_{\lambda})\) by \(\mathrm{E}_{\lambda}\) ; since \(\mathrm{V}_{\lambda}\) is a nonzero \(\mathrm{D}_{\lambda}\) -vector space, it is a simple \(\mathrm{E}_{\lambda}\) -module, (VIII, p.~45, Example 3), and its commutant is equal to \((\mathrm{D}_{\lambda})_{\mathrm{V}_{\lambda}}\) (VIII, p.~82, Corollary 1 of Theorem 2). Since \(\mathrm{E}_{\lambda}\) is the ring of homotheties of the B-module \(\mathrm{V}_{\lambda}\) (VIII, p.~71, Corollary of Proposition 9), this proves b).

The canonical description \(\alpha_{M}\) of M defines an isomorphism \(\alpha_{\lambda}\) from \(V_{\lambda}\otimes_{D_{\lambda}^{\circ}}S_{\lambda}\) to \(M_{\lambda}\) . Since \(V_{\lambda}\) is a simple B-module, the B-module \(V_{\lambda}\otimes_{D_{\lambda}^{\circ}}S_{\lambda}\) is isotypical of type \(V_{\lambda}\) (VIII, p.~61, Proposition 1); the same therefore holds for the B-module \(M_{\lambda}\) , which proves a).

By Remark 1 above, there exist elements \(e_{\lambda}\) of B, for \(\lambda\) running through \(\mathcal{S}_{\mathrm{M}}\) , such that \((e_{\lambda})_{\mathrm{V}_{\lambda}} = 1_{\mathrm{V}_{\lambda}}\) and \((e_{\lambda})_{\mathrm{V}_{\mu}} = 0\) for \(\mu \in \mathcal{S}_{\mathrm{M}}\) with \(\mu \neq \lambda\) . The simple B-modules \(V_{\lambda}\) are therefore pairwise nonisomorphic, which proves c) and the first assertion of d). The B-module \(M_{\lambda}\) is isomorphic to \(V_{\lambda} \otimes_{D_{\lambda}^{\circ}} S_{\lambda}\) , so \(\dim_{D_{\lambda}}(S_{\lambda})\) is the multiplicity of \(V_{\lambda}\) in M (II, §3, No.~7, p.~255, Corollary 1).

The mapping \(\sum_{\lambda\in\mathcal{S}_{M}}\alpha_{\lambda}\) from \(\bigoplus_{\lambda\in\mathcal{S}_{M}}V_{\lambda}\otimes_{D_{\lambda}^{\circ}}S_{\lambda}\) to M provides a description (VIII, p.~69, Definition 5) of the semisimple B-module M. By VIII, p.~70, Proposition 8, b), for every \(\lambda\in\mathcal{S}_{M}\) , the mapping from \(S_{\lambda}\) to \(\mathrm{Hom}_{\mathrm{B}}(\mathrm{V}_{\lambda},\mathrm{M})\) described in e) is bijective and \(D_{\lambda}\) -linear. Since it is obviously A-linear, this proves e).

Remark 2. --- It follows from the proof that the mapping

\[
\sum_ {\lambda \in \mathcal {S} _ {\mathrm{M}}} V _ {\lambda} \otimes_ {D _ {\lambda} ^ {\mathrm{o}}} S _ {\lambda} \longrightarrow M
\]

induced by the canonical description of M is a description of the countermodule of M.

PROPOSITION 6. --- a) Viewed as an (A, B°)-bimodule, M is semisimple.

\begin{enumerate}
\def\labelenumi{\alph{enumi})}
\setcounter{enumi}{1}
\item
  For every \(\lambda \in \mathcal{S}_{\mathrm{M}}\) , \(\mathrm{M}_{\lambda}\) is a simple (A, B°)-sub-bimodule of M.
\item
  For every \((\mathrm{A},\mathrm{B}^{\circ})\) -sub-bimodule N of M, there exists a unique subset \(\Lambda\) of \(\mathcal{S}_{\mathrm{M}}\) such that N is equal to \(\oplus_{\lambda \in \Lambda}\mathrm{M}_{\lambda}\) .
\end{enumerate}

Let \(\lambda\) be in \(\mathcal{S}_{\mathrm{M}}\) . The left A-module \(S_{\lambda}\) and the right B-module \(V_{\lambda}\) are simple, and \(D_{\lambda}\) is the opposite ring of the field of endomorphisms of \(S_{\lambda}\) . By Corollary 2 of VIII, p.~63, the (A, B°)-bimodule \(S_{\lambda} \otimes_{D_{\lambda}} V_{\lambda}\) is simple, and the same holds for \(M_{\lambda}\) , which is isomorphic to it. This proves b), and a) follows.

If \(\lambda\) and \(\mu\) are distinct in \(S_{M}\) , then \(M_{\lambda}\) and \(M_{\mu}\) are not isomorphic as A-modules, nor a fortiori as \((A,B^{\circ})\) -bimodules. Assertion c) follows by Corollary 2 of VIII, p.~68.

PROPOSITION 7. --- a) For every element c of the bicommutant C of M and every \(\lambda\in\mathcal{S}_{M}\) , there exists a unique element \(c_{\lambda}\) of \(\operatorname{End}_{\mathrm{D}_{\lambda}}(\mathrm{S}_{\lambda})\) such that for every \(\varphi\in\operatorname{Hom}_{\mathrm{A}}(\mathrm{S}_{\lambda},\mathrm{M})\) and every \(s\in S_{\lambda}\) , we have \(c\varphi(s)=\varphi(c_{\lambda}s)\) .

\begin{enumerate}
\def\labelenumi{\alph{enumi})}
\setcounter{enumi}{1}
\item
  Endow the \(S_{\lambda}\) , for \(\lambda\) running through \(S_{M}\) , with the C-module structure defined by a). Then the canonical mapping \(\alpha_{M}\) from \(\bigoplus_{\lambda\in S_{M}}S_{\lambda}\otimes_{D_{\lambda}}V_{\lambda}\) to M is an isomorphism of \((\mathrm{C},\mathrm{B}^{\circ})\) -bimodules.
\item
  The mapping \(c \mapsto (c_{\lambda})_{\lambda \in \mathcal{S}_{\mathrm{M}}}\) is an isomorphism from C to \(\prod_{\lambda \in \mathcal{S}_{\mathrm{M}}}\mathrm{End}_{\mathrm{D}_{\lambda}}(\mathrm{S}_{\lambda})\) .
\end{enumerate}

Assertions a) and c) follow from VIII, p.~71, Proposition 9 because the canonical mapping \(\alpha_{\mathrm{M}}\) from \(\bigoplus_{\lambda \in \mathcal{S}} (\mathrm{S}_{\lambda} \otimes_{\mathrm{D}_{\lambda}} \mathrm{W}_{\lambda})\) to M provides a description of the B-module M (Remark 2). Moreover, \(\alpha_{\mathrm{M}}\) is (C, B°)-linear, which proves b).

Remark 3. --- Endow \(S_{\lambda}\) , for \(\lambda \in S_{M}\) , with the C-module structure given by Proposition 7, a). If we replace A with B and B with C in Proposition 5 (VIII, p.~85), then we see that for every \(\lambda \in S_{M}\) , the left C-module \(S_{\lambda}\) is simple, with commutant \(D_{\lambda}\) , that the isotypical component of type \(S_{\lambda}\) of the C-module M is equal to \(M_{\lambda}\) , and that the mapping \(\lambda \mapsto \mathrm{cl}_{\mathrm{C}}(\mathrm{S}_{\lambda})\) is a bijection from the support of the A-module M to the support of the C-module M. Finally, observe that the A-linear mappings and C-linear mappings from \(S_{\lambda}\) to M are identical, as are the A-submodules and C-submodules of M, and that the rings \(\operatorname{End}_{\mathrm{A}}(\mathrm{M})\) and \(\operatorname{End}_{\mathrm{C}}(\mathrm{M})\) are equal.

Let M be a semisimple A-module. Denote by Z the center of the bi-commutant C of the A-module M; it is also the center of the commutant B of M. Endow M and the \(S_{\lambda}\) , for \(\lambda \in S_{M}\) , with Z-module structures deduced by restriction of scalars from the C-module structures. For every \(\lambda \in S_{M}\) , denote the center of the field \(D_{\lambda}\) by \(Z_{\lambda}\) .

PROPOSITION 8. --- a) The mapping \(z \mapsto (z_{\mathrm{S}_{\lambda}})_{\lambda \in \mathcal{S}_{\mathrm{M}}}\) is an isomorphism from Z to the product of the fields \(\mathbf{Z}_{\lambda}\) .

\begin{enumerate}
\def\labelenumi{\alph{enumi})}
\setcounter{enumi}{1}
\item
  The A-module M is isotypical and nonzero if and only if Z is a field.
\item
  Let \(\Lambda\) be a subset of \(\mathcal{S}_{\mathrm{M}}\) . Denote by \(e_{\Lambda}\) the unique element of \(\mathrm{Z}\) such that \((e_{\Lambda})_{\mathrm{S}_{\lambda}} = 1_{\mathrm{S}_{\lambda}}\) for \(\lambda \in \Lambda\) and \((e_{\Lambda})_{\mathrm{S}_{\lambda}} = 0\) for \(\lambda \in \mathcal{S}_{\mathrm{M}} - \Lambda\) . We have \((e_{\Lambda})_{\mathrm{M}_{\lambda}} = 1_{\mathrm{M}_{\lambda}}\) for \(\lambda \in \Lambda\) and \((e_{\Lambda})_{\mathrm{M}_{\lambda}} = 0\) for \(\lambda \in \mathcal{S}_{\mathrm{M}} - \Lambda\) .
\item
  If the support \(S_{M}\) of M is finite, then the mapping \(\Lambda \mapsto e_{\Lambda}Z\) is a bijection from the set of subsets of \(S_{M}\) to the set of ideals of Z, and the mapping \(a \mapsto aM\) is a bijection from the set of ideals of Z to the set of \((A,B^{o})\) -sub-bimodules of M. These bijections are isomorphisms of ordered sets. The reverse bijection sends an \((A,B^{o})\) -sub-bimodule N of M to the ideal consisting of the elements z of Z that send N into M.
\end{enumerate}

For \(\lambda \in \mathcal{S}_{\mathrm{M}}\) , \(Z_{\lambda}\) is the common center of the commutant \(D_{\lambda}\) and the bi-commutant \(C_{\lambda}\) of the A-module \(S_{\lambda}\) . By Proposition 7, c) above, the mapping \(c \mapsto (c_{\lambda})_{\lambda \in \mathcal{S}_{\mathrm{M}}}\) is an isomorphism from C to \(\prod_{\lambda \in \mathcal{S}_{\mathrm{M}}} C_{\lambda}\) . By restriction to the centers, we obtain the isomorphism \(z \mapsto (z_{S_{\lambda}})_{\lambda \in \mathcal{S}_{\mathrm{M}}}\) from Z to \(\prod_{\lambda \in \mathcal{S}_{\mathrm{M}}} Z_{\lambda}\) , whence a).

The ring \(\prod_{\lambda \in \mathcal{S}_{\mathrm{M}}}Z_{\lambda}\) is a field if and only if the set \(\mathcal{S}_{\mathrm{M}}\) has a single element, whence b).

Assertion c) follows from Proposition 7, a). Suppose that \(\mathcal{S}_{\mathrm{M}}\) is finite. It follows from a) and Proposition 8 of I, §8, No.~10, p.~109, that the mapping \(\Lambda \mapsto e_{\Lambda}\mathrm{Z}\) is an isomorphism of ordered sets from \(\mathfrak{P}(\mathcal{S}_{\mathrm{M}})\) to the set of ideals of Z. Let \(\Lambda\) be a subset of \(\mathcal{S}_{\mathrm{M}}\) . By c), we have the relation \(e_{\Lambda}\mathrm{ZM} = e_{\Lambda}\mathrm{M} = \bigoplus_{\lambda \in \Lambda}\mathrm{M}_{\lambda}\) ; because of Proposition 6, c) of VIII, p.~86, it remains to describe the inverse bijection. But \(z \in \mathrm{Z}\) sends M into \(\bigoplus_{\lambda \in \Lambda}\mathrm{M}_{\lambda}\) if and only if \(z = e_{\Lambda}z\) , that is, \(z \in e_{\Lambda}\mathrm{Z}\) .

COROLLARY. --- Suppose that A is an algebra over an algebraically closed commutative field K and that M is a semisimple A-module that is finite-dimensional as a vector space over K. For every \(\lambda\) in \(S_{M}\) , denote by \(e_{\lambda}\) the projector of M with image \(M_{\lambda}\) and kernel \(\oplus_{\lambda\neq\mu}M_{\mu}\) . Then \((e_{\lambda})_{\lambda\in\mathcal{S}_{M}}\) is a basis of the vector space Z over K.

Since M is a finite-dimensional vector space over K that is the direct sum of the family of nonzero submodules \((\mathrm{M}_{\lambda})_{\lambda\in\mathcal{S}_{\mathrm{M}}}\) , the set \(S_{M}\) is finite, and each of the spaces \(S_{\lambda}\) , for \(\lambda\in S_{M}\) , is finite-dimensional over K. Since the field K is algebraically closed, we have \(D_{\lambda}=Z_{\lambda}=K\) (VIII, p.~47, Theorem 1), and the mapping \(z\mapsto(z_{\mathrm{S}_{\lambda}})_{\lambda\in\mathcal{S}_{\mathrm{M}}}\) is an isomorphism from Z to \(K^{S_{M}}\) (Proposition 8, a)). The corollary then follows from part c) of Proposition 8.

\subsection{Density Theorem}

THEOREM 3 (Jacobson). --- Let M be a semisimple A-module, and let c be an endomorphism of the additive group M. Then c belongs to the bicommutant \(A_{M}^{\prime\prime}\) of M if and only if it satisfies the following condition:

\begin{enumerate}
\def\labelenumi{(\Alph{enumi})}
\setcounter{enumi}{3}
\tightlist
\item
  For every finite subset \(\mathrm{F}\) of \(\mathrm{M}\) , there exists an element \(a\) of \(\mathrm{A}\) such that \(c\) coincides with \(a_{\mathrm{M}}\) on \(\mathrm{F}\) .
\end{enumerate}

First, suppose that \(c\) satisfies condition (D). Let \(u\) be an element of \(\mathrm{A}_{\mathrm{M}}^{\prime}\) . Let \(x\) be an element of \(\mathrm{M}\) , and apply condition (D) to the subset \(\mathrm{F} = \{x, u(x)\}\) . There exists an element \(a\) of \(\mathrm{A}\) such that \(c(x) = ax\) and \(c(u(x)) = au(x)\) , so that \(u(c(x)) = u(ax) = au(x) = c(u(x))\) . Since \(x\) is arbitrary, we have \(cu = uc\) ; as this holds for all \(u\) , we have \(c \in \mathrm{A}_{\mathrm{M}}^{\prime \prime}\) .

For the converse, we use the following lemma.

Lemma 3. --- Let M be a semisimple A-module. Let B be the bicommutant of the A-module M. Then every A-submodule of M is a B-submodule of M.

Let N be an A-submodule of M. By Corollary 2 of VIII, p.~56, there exists a projector p of the A-module M with image N; it belongs to \(A_{M}^{\prime}\) . Since we have the relation pb = bp for every \(b \in B\) , we obtain that N is a B-submodule of M.

Let us conclude the proof of Theorem 3. Suppose that c belongs to \(A_{M}^{\prime\prime}\) , and let \(F = \{x_{1}, \ldots, x_{n}\}\) be a finite subset of M. Denote the element \((x_{1}, \ldots, x_{n})\) of \(M^{n}\) by x. The A-module \(M^{n}\) is semisimple and, by Proposition 2 of VIII, p.~79, its bicommutant coincides with the homotheties of the \(A_{M}^{\prime\prime}\) -module \(M^{n}\) . By Lemma 3, the A-submodule Ax of \(M^{n}\) is an \(A_{M}^{\prime\prime}\) -submodule of \(M^{n}\) . Let \(a \in A\) be such that \((cx_{1}, \ldots, cx_{n})\) is equal to ax. Then c coincides with \(a_{M}\) on \(\{x_{1}, \ldots, x_{n}\}\) , which implies condition (D).

Remark. --- Denote the endomorphism ring of the additive group of M by E. Endow M with the discrete topology; the ring E consists of mappings from M to M, and we can endow it with the topology induced by the product topology on \(M^{M}\) (``topology of simple convergence in M'', Gen.~Top., I, §2, No.~3, p.~31). The topology on E is Hausdorff and compatible with the additive group structure on E. For every f in E, the mappings \(g \mapsto f \circ g\) and \(g \mapsto g \circ f\) from E to E are continuous. Consequently, the commutant of every subset of E is closed in E. Theorem 3 therefore implies that \(A_{M}^{\prime\prime}\) is the closure of \(A_{M}\) in E.

\subsection{Application to Field Theory}

PROPOSITION 9. --- Let L be a field and E a subring of \(\operatorname{End}_{\mathbf{Z}}(\mathrm{L})\) that contains the mapping \(\gamma_{a}: x \mapsto ax\) for every \(a \in L\) . Denote by K the set of elements a of L such that \(u(xa) = u(x)a\) for every x in L and every u in E; it is a subfield of L.

Let V be a finite-dimensional linear subspace of the right K-vector space L, and let h be a K-linear mapping from V to L. There exists an element of E that coincides with h on V.

We view L as a left E-module. Since E contains the left multiplications \(\gamma_{a}\) , every E-submodule of L is a left ideal of the field L; the E-module L is therefore simple. Every endomorphism of the additive group of L that commutes with the \(\gamma_{a}\) is of the form \(\delta_{b}: x \mapsto xb\) with b in L. Consequently, \(b \mapsto \delta_{b}\) is an isomorphism from K to the opposite ring of \(\operatorname{End}_{\operatorname{E}}(\operatorname{L})\) , which is a field.

The bicommutant \(E''\) of the E-module L therefore consists of the endomorphisms of the right K-vector space L. Let v be an endomorphism of the K-vector space L whose restriction to V coincides with h; it is an element of \(E''\) . Let \((x_{i})_{i\in\mathrm{I}}\) be a basis of V over K; by Theorem 3 (VIII, p.~88), there exists an element u of E such that \(u(x_{i}) = v(x_{i}) = h(x_{i})\) for \(i \in I\) . By linearity, it follows that \(u(x) = h(x)\) for every x in V.

COROLLARY. --- Let L be a field. Let \(\Gamma\) be a subgroup of the automorphism group of the field L, and let K be the field of invariants of \(\Gamma\) . Let V be a right K-linear subspace of L of finite dimension n over K. Then there exist elements \(\sigma_{1},\ldots,\sigma_{n}\) of \(\Gamma\) with the following property: for every K-linear mapping u from V to L, there exist elements \(a_{1},\ldots,a_{n}\) of L such that we have \(u(x)=\sum_{i=1}^{n}a_{i}\sigma_{i}(x)\) for every x in V.

Denote by E the set of mappings from L to L of the form \(x \mapsto \Sigma_{\sigma \in \Gamma} a_{\sigma} \sigma(x)\) , where \((a_{\sigma})_{\sigma \in \Gamma}\) is a family of elements of L with finite support. We have \(\gamma_{a} \in E\) for every a in L, and E is a subring of the endomorphism ring of the additive group of L. Moreover, the field K consists of the elements a of L such that \(u(xa) = u(x)a\) for every x in L and every u in E.

Let H be the left L-vector space \(\operatorname{Hom}_{\mathrm{K}}(\mathrm{V},\mathrm{L})\) ; it has dimension n.~By Proposition 9, it is generated by the restrictions of the elements of \(\Gamma\) to V. There exist n elements \(\sigma_{1},\ldots,\sigma_{n}\) of \(\Gamma\) whose restrictions to V form a basis of H over L. The corollary follows from this.

Remark. --- When the field L is commutative, this corollary reduces to Artin's theorem (V, §10, No.~6, p.~65).

\BourbakiExercises{Exercises}

\begin{enumerate}
\def\labelenumi{\arabic{enumi})}
\tightlist
\item
  Let M be an A-module and N a direct factor submodule of M.
\end{enumerate}

\begin{enumerate}
\def\labelenumi{\alph{enumi})}
\item
  Prove that N is stable under the bicommutant \(A_{M}^{\prime\prime}\) of M.
\item
  Prove that every A-linear mapping from N to M is \(A_{M}^{\prime\prime}\) -linear.
\item
  Prove that restriction to N defines a ring homomorphism \(A_{M}^{\prime\prime} \to A_{N}^{\prime\prime}\) .
\item
  Suppose that M/N is a sum of submodules isomorphic to quotients of N. Prove that the canonical homomorphism \(A_{M}^{\prime\prime} \rightarrow A_{N}^{\prime\prime}\) is injective. If N is balanced, the same holds for M.
\end{enumerate}

\begin{enumerate}
\def\labelenumi{\arabic{enumi})}
\setcounter{enumi}{1}
\tightlist
\item
  Let K be a commutative field and A the subring of \(\mathbf{M}_{3}(\mathrm{K})\) consisting of the matrices of the form
\end{enumerate}

\[
\left( \begin{array}{c c c} a & 0 & 0 \\ b & c & 0 \\ 0 & 0 & a \end{array} \right)
\]

for \(a, b, c\) in K. Let M be the A-module \(\mathrm{K}^3\) ; it is the direct sum of the submodules \(\mathrm{N} = \mathrm{Ke}_1 + \mathrm{Ke}_2\) and \(\mathrm{P} = \mathrm{Ke}_3\) .

\begin{enumerate}
\def\labelenumi{\alph{enumi})}
\item
  Prove that M is balanced, that the canonical mapping \(A_{M}^{\prime\prime} \rightarrow A_{N}^{\prime\prime}\) (Exercise 1) is injective but not surjective, and that the canonical mapping \(A_{M}^{\prime\prime} \rightarrow A_{P}^{\prime\prime}\) is surjective but not injective.
\item
  Deduce an example of a module Q over a ring B and a direct factor submodule R such that the canonical mapping \(B_{Q}^{\prime\prime} \rightarrow B_{R}^{\prime\prime}\) is neither injective nor surjective (take \(B = A \times A\) and \(Q = M \times M\) ).
\end{enumerate}

\begin{enumerate}
\def\labelenumi{\arabic{enumi})}
\setcounter{enumi}{2}
\item
  Let B be a ring and I a two-sided ideal of B; let A be the subring of \(B \times B\) consisting of the pairs \((x, y)\) such that \(x - y \in I\) . Let \(B_{1}\) (resp. \(B_{2}\) ) be the additive group B endowed with an A-module structure deduced from the first (resp. second) projection. Prove that the A-modules \(B_{1}\) and \(B_{2}\) are balanced. Under what condition on I is the A-module \(B_{1} \oplus B_{2}\) balanced?
\item
  Let A and B be rings and P an (A, B)-bimodule.
\end{enumerate}

\begin{enumerate}
\def\labelenumi{\alph{enumi})}
\item
  Suppose that the B-module P is generating and that the A-module P is balanced. Show that for every balanced B-module M, the A-module \(P \otimes_{B} M\) is balanced.
\item
  Suppose that the A-module P is generating and that the mapping \(b \mapsto b_{P}\) from B to \(\operatorname{End}_{\mathrm{A}}(\mathrm{P})\) is bijective. Prove that for every generating B-module M, the A-module \(P \otimes_{B} M\) is generating.
\end{enumerate}

\begin{enumerate}
\def\labelenumi{\arabic{enumi})}
\setcounter{enumi}{4}
\item
  Let A be a ring, M an A-module, and N a submodule of M. Prove that if the A-module M/N is generating, the same holds for the A-module M. Use an example to prove that N can be generating without M being so.
\item
  Let \((\mathrm{A}_i)_{i\in \mathrm{I}}\) be a family of rings and A be its product. Prove that the A-module \(\bigoplus \mathrm{A}_i\) is projective and faithful. Is it always generating?
\item
  Let A be a ring and M an A-module. Prove that if M is generating, the right A-module \(\mathrm{M}^{*} = \mathrm{Hom}_{\mathrm{A}}(\mathrm{M}, \mathrm{A}_{s})\) is generating. Use an example to prove that \(M^{*}\) can be generating without M being so (cf.~Exercise 6).
\item
  Let A be a principal ideal domain and M an A-module. Prove that the following properties are equivalent:
\end{enumerate}

\begin{enumerate}
\def\labelenumi{(\roman{enumi})}
\item
  The A-module M is generating.
\item
  The dual of M is nonzero.
\item
  The module M contains a direct factor submodule isomorphic to \(A_{s}\) .
\end{enumerate}

\begin{enumerate}
\def\labelenumi{\arabic{enumi})}
\setcounter{enumi}{8}
\tightlist
\item
  Let \(D\) be a field and \(V\) a right vector space over \(D\) . We view \(V\) as a left module over the ring \(A = \operatorname{End}_D(V)\) .
\end{enumerate}

\begin{enumerate}
\def\labelenumi{\alph{enumi})}
\item
  Prove that the A-module V is projective, finitely generated, faithful, and balanced. Prove that the (D,A)-bimodules \(\mathrm{Hom}_{\mathrm{A}}(\mathrm{V},\mathrm{A}_{s})\) and \(\mathrm{V}^{*}=\mathrm{Hom}_{\mathrm{D}}(\mathrm{V},\mathrm{D}_{d})\) are canonically isomorphic.
\item
  Prove that the A-module V is generating if and only if the dimension of the D-vector space V is finite.
\item
  What is the trace ideal of the A-module V?
\end{enumerate}

\begin{enumerate}
\def\labelenumi{\arabic{enumi})}
\setcounter{enumi}{9}
\tightlist
\item
  Let M be an A-module. Prove that the following properties are equivalent:
\end{enumerate}

\begin{enumerate}
\def\labelenumi{(\roman{enumi})}
\item
  The A-module M is finitely generated.
\item
  For every generating A-module P, there exist a natural number n and a surjective homomorphism f from \(P^{n}\) to M.
\end{enumerate}

\begin{enumerate}
\def\labelenumi{\arabic{enumi})}
\setcounter{enumi}{10}
\tightlist
\item
  Let \(\mathrm{P}\) be a finitely generated projective A-module; denote its trace ideal by \(\tau(\mathrm{P})\) .\\
\end{enumerate}

\begin{enumerate}
\def\labelenumi{\alph{enumi})}
\item
  Prove that we have \(\tau(\mathrm{P})^2 = \tau(\mathrm{P})\) and \(\tau(\mathrm{P})\mathrm{P} = \mathrm{P}\) .
\item
  Prove that the canonical mapping \(\tau(\mathrm{P})\otimes_{\mathrm{A}}\mathrm{P}\to\mathrm{P}\) is an isomorphism.
\end{enumerate}

¶ 12) Let A be a ring, and let P be a finitely generated A-module. Prove that the following properties are equivalent:

\begin{enumerate}
\def\labelenumi{(\roman{enumi})}
\item
  The A-module P is projective and generating.
\item
  The A-modules \(\mathrm{P}^{(\mathbf{N})}\) and \(\mathrm{A}_s^{(\mathbf{N})}\) are isomorphic.
\end{enumerate}

(To establish the implication \((\mathrm{i})\Rightarrow(\mathrm{ii})\) , use induction to construct strictly increasing sequences \((n_{k})\) and \((m_{k})\) of integers and surjective homomorphisms of A-modules \(f_{k}:A_{s}^{n_{k}}\to P^{m_{k}}\) and \(g_{k}:P^{m_{k+1}}\to A_{s}^{n_{k}}\) such that \(g_{k-1}\circ f_{k}\) and \(f_{k}\circ g_{k}\) are the canonical projections.)

\begin{enumerate}
\def\labelenumi{\arabic{enumi})}
\setcounter{enumi}{12}
\item
  Let k be a ring and n an integer \(\geqslant 2\) ; denote by A the subring of \(\mathbf{M}_{n}(k)\) consisting of the upper-triangular matrices and by M the right A-module \(k^{n}\) . Prove that M is projective and finitely generated but that the ring of homotheties of M is not dense in its bicommutant (cf.~VIII, p.~88, Theorem 3 and VIII, p.~88, Remark). In particular, M is not balanced.
\item
  What is the bicommutant of the Z-module Q? Prove that Q is a simple module over its bicommutant but that the ring of homotheties of the Z-module Q is not dense in its bicommutant.
\end{enumerate}

¶ 15) Let M be a torsion module over a principal ideal domain. Prove that the ring of homotheties of M is dense in its bicommutant.

(Reduce to the case when M is \(\pi\) -primary for an irreducible element \(\pi\) of A. Then, prove that every finite subset S of M is contained in a direct factor submodule that is the direct sum of a finitely generated module and a finite number of indecomposable divisible modules. Use induction on the cardinal of S with the help of Exercises 3 of VII, §2, p.~54 and 8 of VII, §2, p.~56.)

\begin{enumerate}
\def\labelenumi{\arabic{enumi})}
\setcounter{enumi}{15}
\tightlist
\item
  Let D be a field, V a vector space of dimension \(\geqslant 2\) over D, and A a subring of \(\operatorname{End}_{\mathrm{D}}(\mathrm{V})\) .
\end{enumerate}

\begin{enumerate}
\def\labelenumi{\alph{enumi})}
\item
  Suppose that for every system of four elements \((x, x', y, y')\) , where x and \(x'\) are linearly independent, there exists an element u of A such that \(u(x) = y\) and \(u(x') = y'\) . Prove that the commutant of the A-module V is D (identified with the ring of homotheties of V) and that the ring of homotheties of this module is dense in its bicommutant.
\item
  Give an example of a ring \(A \subset \operatorname{End}_{\mathrm{D}}(\mathrm{V})\) such that V is a simple A-module but that its commutant is distinct from D (consider a field E containing D, a vector space V over E, and the ring \(A = \operatorname{End}_{\mathrm{E}}(\mathrm{V})\) ).
\end{enumerate}

\begin{enumerate}
\def\labelenumi{\arabic{enumi})}
\setcounter{enumi}{16}
\tightlist
\item
  Let K be an algebraically closed commutative field, V a finite-dimensional vector space over K, and G a submonoid of \(\operatorname{Aut}(V)\) such that the K{[}G{]}-module V is simple. Prove that if the set of scalars \(\operatorname{Tr}(g)\) , for \(g \in G\) , is finite, then G is finite (deduce from Burnside's theorem that G contains a basis \((g_{i})_{i \in I}\) of \(\operatorname{End}(V)\) , and consider the mapping \(g \mapsto (\operatorname{Tr}(gg_{i}))\) from G to \(K^{I}\) ).
\end{enumerate}

¶ 18) A group G is called locally finite if every subgroup of G that admits a finite generating family is finite.

\begin{enumerate}
\def\labelenumi{\alph{enumi})}
\item
  Let G be a group and H a normal subgroup of G. Prove that if H and G/H are locally finite, the same holds for G.
\item
  Deduce from a) that a solvable group whose elements are all of finite order is locally finite.
\item
  Let K be a commutative field, V a finite-dimensional vector space over K, and G a subgroup of Aut(V) that admits a finite generating family consisting of elements of finite order. Prove that the order of the elements of G is bounded.
\end{enumerate}

(Reduce to the case when K is a purely transcendental extension of a prime field P by considering the subfield K generated by the coefficients of a finite generating family of G. Then, prove that for every element s of \(\mathbf{GL}_{n}(\mathrm{K})\) of finite order, we have \(\varphi(s) \leqslant n\) if P = Q and \(s \leqslant p^{n} - 1\) if \(P = F_{p}\) .)

\begin{enumerate}
\def\labelenumi{\alph{enumi})}
\setcounter{enumi}{3}
\tightlist
\item
  Prove that every subgroup H of Aut(V) whose elements are all of finite order is locally finite. \(^{(1)}\)
\end{enumerate}

(Reduce to the case when K is algebraically closed, and use induction on \(\dim(V)\) . Let G be a subgroup of H that admits a finite generating family. If V is a simple K{[}G{]}-module, apply Exercise 17; if V contains a nontrivial subspace W that is stable under G, consider the image of G in \(\operatorname{Aut}(V) \times \operatorname{Aut}(V/W)\) , and apply a) and b)).

¶ 19) Let A be a ring such that the A-module \(A_{s}\) is the direct sum of a finite family of left ideals \((\mathfrak{I}_{i})_{1\leqslant i\leqslant n}\) that are isomorphic as A-modules.

\begin{enumerate}
\def\labelenumi{\alph{enumi})}
\item
  Prove that in A, there exists a family of elements \(e_{i,j}\) ( \(1 \leqslant i \leqslant n\) , \(1 \leqslant j \leqslant n\) ) such that \(e_{i,j}e_{h,k} = \delta_{j,h}e_{i,k}\) , where \(\delta_{j,h}\) denotes the Kronecker delta function, and \(I_i = Ae_{i,i}\) (cf.~VIII, p.~39, Exercise 5 and I, §8, p.~173, Exercise 12). Conversely, if in A, there exists a family of \(n^2\) elements \(e_{i,j}\) such that \(e_{i,j}e_{h,k} = \delta_{j,h}e_{i,k}\) and \(1 = \sum_{i=1}^{n} e_{i,i}\) , then the left ideals \(Ae_{i,i}\) are isomorphic as A-modules (cf.~VIII, p.~39, Exercise 5). Moreover, if B is the subring of A consisting of the elements that commute with all the \(e_{i,j}\) , then A is isomorphic to the matrix ring \(\mathbf{M}_n(\mathbf{B})\) and B is isomorphic to the opposite ring of the commutant of each of the A-modules \(Ae_{i,i}\) (use VIII, p.~78 and VIII, p.~65, Corollary).
\item
  Let M be an A-module. Prove that the \(N_{i} = e_{i,i}M\) are B-modules and that M, viewed as a B-module, is the direct sum of the \(N_{i}\) ; moreover, the B-modules \(N_{i}\) are pairwise isomorphic (consider the mappings \(x \mapsto e_{i,1}x\) and \(y \mapsto e_{1,i}y\) ), and the annihilator of the \(N_{i}\) is the intersection of B and the annihilator of M in A. Conversely, for every B-module N, define an A-module structure on the direct sum M of n B-modules isomorphic to N, so that \(e_{1,1}M\) is isomorphic to N.
\item
  To each B-submodule P of \(N_{1}\) , assign the A-submodule \(\sum_{i=1}^{n} e_{i,1}P\) of M. Prove that this defines a bijective mapping from the set of B-submodules of \(N_{1}\) to the set of A-submodules of M that is strictly increasing for the inclusion relation.
\end{enumerate}

\section{Morita Equivalence of Modules and Algebras}

In this section, \(k\) denotes a commutative ring.

\subsection{Commutant and Duality}

Let A and B be k-algebras. Recall (III, §4, No.~3, p.~466) that a bimodule over the algebras A and B is an (A,B)-bimodule P on which the two k-module structures deduced from the A- and B-module structures coincide. To avoid any ambiguity, we say that P is an \((\mathrm{A},\mathrm{B})_{k}\) -bimodule. Let P be an \((\mathrm{A},\mathrm{B})_{k}\) -bimodule. We denote by \(P^{*}\) the dual \(\operatorname{Hom}_{\mathrm{A}}(\mathrm{P},\mathrm{A})\) of the left A-module underlying P. It is a \((\mathrm{B},\mathrm{A})_{k}\) -bimodule (II, §1, No.~14, p.~226); for \(a \in A\) , \(b \in B\) , \(x \in P\) , and \(x^{*} \in P^{*}\) , we have

\[
\langle x, b x ^ {*} a \rangle = \langle x b, x ^ {*} \rangle a.\tag{1}
\]

We denote by \(_{s}A_{d}\) the algebra A viewed as an \((\mathrm{A},\mathrm{A})_{k}\) -bimodule (loc. cit.) and by \(\Lambda:P\otimes_{B}P^{*}\to_{s}A_{d}\) the homomorphism of \((\mathrm{A},\mathrm{A})_{k}\) -bimodules determined by

\[
\Lambda (x \otimes x ^ {*}) = \langle x, x ^ {*} \rangle\tag{2}
\]

for \(x \in \mathrm{P}\) and \(x^{*} \in \mathrm{P}^{*}\) . We denote by \(\widetilde{\mathrm{P}}\) the dual \(\mathrm{Hom}_{\mathrm{B}}(\mathrm{P},\mathrm{B})\) of the right B-module underlying \(\mathrm{P}\) ; it is a \((\mathrm{B},\mathrm{A})_{k}\) -bimodule. We denote by \(\widetilde{\Lambda}: \widetilde{\mathrm{P}} \otimes_{\mathrm{A}} \mathrm{P} \to {}_{s}\mathrm{B}_{d}\) the homomorphism of \((\mathrm{B},\mathrm{B})_{k}\) -bimodules determined by

\[
\widetilde {\Lambda} (\widetilde {x} \otimes x) = \langle \widetilde {x}, x \rangle\tag{3}
\]

for \(x \in P\) and \(\widetilde{x} \in \widetilde{P}\) .

Now, suppose that the mapping \(b \mapsto b_{\mathrm{P}}\) is a bijection from B to \(\operatorname{End}_{\mathrm{A}}(\mathrm{P})\) ; it is then an isomorphism from B to the opposite algebra of \(\operatorname{End}_{\mathrm{A}}(\mathrm{P})\) . The canonical homomorphism of \(\mathbf{Z}\) -modules from \(\mathrm{P}^* \otimes_{\mathrm{A}} \mathrm{P}\) to \(\operatorname{End}_{\mathrm{A}}(\mathrm{P})\) (II, §4, No.~2, p.~271) then determines a homomorphism of \(\mathbf{Z}\) -modules \(\Theta: \mathrm{P}^* \otimes_{\mathrm{A}} \mathrm{P} \to \mathrm{B}\) defined by

\[
x \Theta (x ^ {*} \otimes y) = \langle x, x ^ {*} \rangle y\tag{4}
\]

for \(x, y \in \mathrm{P}\) and \(x^{*} \in \mathrm{P}^{*}\) . Because of (1), this homomorphism is (B,B)-linear, and we have

\[
\Theta (x ^ {*} \otimes y) y ^ {*} = x ^ {*} \langle y, y ^ {*} \rangle\tag{5}
\]

for \(y \in \mathrm{P}\) and \(x^{*}, y^{*} \in \mathrm{P}^{*}\) . From (4) and (5), we deduce the following equalities in the \((\mathrm{B}, \mathrm{B})_{k}\) -bimodule \(\mathrm{P}^{*} \otimes_{\mathrm{A}} \mathrm{P}\) :

\[
(y ^ {*} \otimes x) \Theta (x ^ {*} \otimes y) = y ^ {*} \otimes \langle x, x ^ {*} \rangle y = y ^ {*} \langle x, x ^ {*} \rangle \otimes y = \Theta (y ^ {*} \otimes x) (x ^ {*} \otimes y)
\]

for \(x, y \in \mathrm{P}\) and \(x^{*}, y^{*} \in \mathrm{P}^{*}\) , and therefore

\[
s \Theta (t) = \Theta (s) t\tag{6}
\]

for \(s, t \in \mathrm{P}^* \otimes_{\mathrm{A}} \mathrm{P}\) . Likewise, for \(x, y\) in \(\mathrm{P}\) and \(x^*, y^*\) in \(\mathrm{P}^*\) , we deduce the following equalities in the \((\mathrm{A}, \mathrm{A})_k\) -bimodule \(\mathrm{P} \otimes_{\mathrm{B}} \mathrm{P}^*\) from (4) and (5):

\[
(x \otimes x ^ {*}) \langle y, y ^ {*} \rangle = x \otimes \Theta (x ^ {*} \otimes y) y ^ {*} = x \Theta (x ^ {*} \otimes y) \otimes y ^ {*} = \langle x, x ^ {*} \rangle (y \otimes y ^ {*}),
\]

and therefore

\[
u \Lambda (v) = \Lambda (u) v\tag{7}
\]

for \(u, v \in \mathrm{P} \otimes_{\mathrm{B}} \mathrm{P}^{*}\) .

For any element \(x^{*}\) of \(P^{*}\) , denote the B-linear mapping \(x \mapsto \Theta(x^{*} \otimes x)\) from P to B by \(\sigma(x^{*})\) . We thus define a mapping \(\sigma\) from \(P^{*}\) to \(\widetilde{P}\) that is (B, A)-linear and satisfies, by definition,

\[
\Theta (x ^ {*} \otimes y) = \langle \sigma (x ^ {*}), y \rangle\tag{8}
\]

for \(x^{*} \in P^{*}\) and \(y \in P\) . By the definition of \(\widetilde{\Lambda}\) , we therefore have

\[
\Theta = \widetilde {\Lambda} \circ (\sigma \otimes 1 _ {\mathrm{P}}).\tag{9}
\]

Suppose that the mapping \(a \mapsto a_{\mathrm{P}}\) from A to \(\operatorname{End}_{\mathrm{B}}(\mathrm{P})\) is bijective; it is then an algebra isomorphism. Analogously, we define a homomorphism of \((\mathrm{A}, \mathrm{A})_k\) -bimodules \(\widetilde{\Theta}: \mathrm{P} \otimes_{\mathrm{B}} \widetilde{\mathrm{P}} \to {}_s \mathrm{A}_d\) by setting

\[
\widetilde {\Theta} (x \otimes \widetilde {y}) y = x \langle \widetilde {y}, y \rangle\tag{10}
\]

for \(x, y \in \mathrm{P}\) and \(\widetilde{y} \in \widetilde{\mathrm{P}}\) . We also define a homomorphism of \((\mathrm{B}, \mathrm{A})_k\) -bimodules \(\widetilde{\sigma}: \widetilde{\mathrm{P}} \to \mathrm{P}^*\) by setting

\[
\widetilde {\Theta} (x \otimes \widetilde {y}) = \langle x, \widetilde {\sigma} (\widetilde {y}) \rangle\tag{11}
\]

for \(x\in \mathrm{P},\widetilde{y}\in \widetilde{\mathrm{P}}\) . We have

\[
\widetilde {\Theta} = \Lambda \circ (1 _ {P} \otimes \widetilde {\sigma}).\tag{12}
\]

PROPOSITION 1. --- Suppose that the mappings \(b \mapsto b_{P}\) from B to \(\operatorname{End}_{\mathrm{A}}(\mathrm{P})\) and \(a \mapsto a_{P}\) from A to \(\operatorname{End}_{\mathrm{B}}(\mathrm{P})\) are bijective. Then \(\sigma\) and \(\widetilde{\sigma}\) are inverse isomorphisms, and we have the relations

\[
\Lambda = \widetilde {\Theta} \circ (1 _ {\mathrm{P}} \otimes \sigma),\tag{13}
\]

\[
\widetilde {\Lambda} = \Theta \circ (\widetilde {\sigma} \otimes 1 _ {\mathrm{P}}).\tag{14}
\]

For \(x \in \mathrm{P}\) , \(x^{*} \in \mathrm{P}^{*}\) , and \(y \in \mathrm{P}\) , by relations (4), (8), (10), and (11), we have

\[
\langle x, x ^ {*} \rangle y = x \Theta (x ^ {*} \otimes y) = x \langle \sigma (x ^ {*}), y \rangle = \widetilde {\Theta} (x \otimes \sigma (x ^ {*})) y = \langle x, \widetilde {\sigma} (\sigma (x ^ {*})) \rangle y.\tag{15}
\]

Likewise, for \(x \in P\) , \(\widetilde{y} \in \widetilde{P}\) , and \(y \in P\) , we have

\[
x \langle \widetilde {y}, y \rangle = \widetilde {\Theta} (x \otimes \widetilde {y}) y = \langle x, \widetilde {\sigma} (\widetilde {y}) \rangle y = x \Theta (\widetilde {\sigma} (\widetilde {y}) \otimes y) = x \langle \sigma (\widetilde {\sigma} (\widetilde {y})), y \rangle .\tag{16}
\]

Because of the assumptions, P is faithful as an A-module and as a B-module. From relations (15) and (16), respectively, we deduce \(\widetilde{\sigma} \circ \sigma = 1_{P^{*}}\) and \(\sigma \circ \widetilde{\sigma} = 1_{\widetilde{P}}\) . Relations (13) and (14) then follow from (12) and (9), respectively.

Remarks. --- 1) Suppose that the mapping \(b \mapsto b_{\mathrm{P}}\) from B to \(\operatorname{End}_{\mathrm{A}}(\mathrm{P})\) is bijective. Then

\begin{enumerate}
\def\labelenumi{\alph{enumi})}
\item
  the B-module P can be identified with the countermodule of P; it is therefore faithful and balanced;
\item
  the mapping \(a \mapsto a_{P}\) from A to \(\operatorname{End}_{\mathrm{B}}(\mathrm{P})\) is bijective if and only if the A-module P is faithful and balanced.
\end{enumerate}

\begin{enumerate}
\def\labelenumi{\arabic{enumi})}
\setcounter{enumi}{1}
\tightlist
\item
  Under the assumptions of Proposition 1, the A-module P is balanced; since the rings \(A_{P}\) and \(A_{P}^{\prime}\) have the same center (VIII, p.~77), there is an isomorphism \(\varphi\) from the center \(Z(A)\) of the ring A to the center \(Z(B)\) of the ring B, determined by the relation \(\varphi(z)_{\mathrm{P}} = z_{\mathrm{P}}\) for \(z \in \mathrm{Z}(A)\) . Moreover, the endomorphisms of the \((\mathrm{A}, \mathrm{B})_{k}\) -bimodule P are the homotheties \(z_{P}\) where z runs through \(Z(A)\) ; the automorphisms of the \((\mathrm{A}, \mathrm{B})_{k}\) -bimodule P are the homotheties \(z_{P}\) where z is invertible in \(Z(A)\) .
\end{enumerate}
